% !TeX root = ../../Book1.tex
% 纲要标题：## 第18章　主丛、规范变换与平坦结构
% 纲要位置：几何学计划纲要.md，第 807 行。
% 正文状态：初稿；尚未正式审核。

% 纲要内容开始
% 
% > 本章目标与证明范围：将过渡函数、水平分布和局部规范场统一，证明主丛联络与伴随丛联络的关系。规范等价和基空间坐标变化必须各有清楚作用对象。
% 
% 纲要内容结束

\chapter{主丛、规范变换与平坦结构}
\label{g1:ch:18}
\BookChapterNumber{g1:ch:18}

\begin{chapterabstract}
	主丛把同一纤维中的标架选择组织为群的自由传递作用。过渡函数描述这些选择如何粘合，主联络则规定沿底空间移动时如何比较它们。本章构造主丛、关联丛与结构群约化，证明水平分布、联络形式和局部联络数据的等价性，并建立关联向量丛上的协变微分。水平提升的全区间存在、曲率结构方程、Bianchi 恒等式及运输变分说明局部微分怎样控制路径比较。最后区分规范变换、换局部截面与底空间坐标变化，以单值表示分类平坦主丛，并用圆周、环面和 Hopf 丛检验这份分类的范围。
\end{chapterabstract}

\begin{learninggoals}
	\begin{itemize}
		\item 从局部截面求过渡函数，由余循环重建主丛，构造关联丛与标架丛，并证明结构群约化的截面判据。
		\item 在水平分布、主联络形式与局部形式之间转换，证明主联络存在，计算联络之差及关联丛上的诱导协变微分。
		\item 证明紧参数区间上路径的水平提升存在且唯一，建立曲率结构方程与 Bianchi 恒等式，用曲率计算运输变分和小环路的首项。
		\item 区分改变局部截面、施行规范变换与改变底空间坐标，计算规范轨道，并用环路运输的中心化子确定规范稳定子。
		\item 由单值表示构造平坦主丛，证明表示模共轭的分类及固定底层丛时的限制，计算圆周、环面的相位和 Hopf 丛的曲率积分。
	\end{itemize}
\end{learninggoals}

给一个向量丛选取局部标架之后，纤维中的向量就有了列坐标。但我们也可以把标架本身当作研究对象：同一点的两组标架之间有唯一可逆矩阵，却没有哪一组天然应当充当单位标架。若要求标架正交，允许的变换便缩为正交群；若所记录的只是一个单位复数相位，允许的变换则是圆周群。主丛用同一种结构容纳这些选择，联络在这个选择空间中规定如何随底点移动。

\ProseParagraphBreak

局部表达与整体比较提出不同的问题。一个变化的标架能使原本为零的联络矩阵出现非零项，但没有改变联络；一条环路的水平运输也可能不平凡，即使曲率处处为零。我们需要同时保留总空间上的几何对象、截面中的表达和环路留下的群元素。平坦结构的分类会把最后一类信息组织成基本群表示，同时说明有些底层主丛根本不容许平坦联络。

\ProseParagraphBreak

本章的 Lie 群均为有限维实 Lie 群，流形默认有限维、无边界、Hausdorff、第二可数；复线性表示另作说明。主丛统一使用右作用 \(p\mapsto pg\)，局部截面满足 \(s_j=s_i g_{ij}\)。因此其基本场取 \(X_p^\#=\left.\frac{\mathrm{d}}{\mathrm{d}t}\right|_0p\exp(tX)\)，括号符号须与前章默认的左作用区分。局部联络公式延续\chapref{g1:ch:15}的行标架与列分量约定；Lie 群及光滑商使用\chapref{g1:ch:16}和\chapref{g1:ch:17}。水平提升、平坦性和单值构造分别调用前面已经建立的常微分方程、Frobenius 定理与万有覆盖，使用时会恢复所需条件。

\ProseParagraphBreak

对照阅读可查 Tu 的主丛、主联络、水平提升、曲率及关联丛各节\cite[第27--31节]{Tu2017Connections}。比较局部公式时，先核对主作用在左还是在右，以及过渡函数从哪一张截面换到哪一张；这些选择决定逆元与共轭的位置。

% !TeX root = ../../Book1.tex
% 纲要标题：### 18.1 主丛和关联构造
% 纲要位置：几何学计划纲要.md，第 811 行。
% 正文状态：初稿；尚未正式审核。

% 纲要内容开始
% 
% * 主作用、局部截面与过渡函数
% * 由余循环重建主丛
% * 关联向量丛和标架丛
% * 结构群约化及截面判据
% 
% 纲要内容结束

\section{主丛和关联构造}
\label{g1:sec:1801}

在向量丛的同一条纤维中，两组标架之间有唯一的可逆换基矩阵，但没有一组标架天然应当充当单位元。主丛保留的正是这种“任意两点可以唯一比较，却未指定起点”的结构。\cref{g1:def:frame-bundle}已经给出了一个实例：标架 \(u:\mathbb R^r\to E_x\) 可以右乘矩阵，\(u\cdot A=u\circ A\)。下面把矩阵群换成一般 Lie 群，先确定局部选择如何相容，再由这些选择构造向量丛及额外几何结构。本节的流形均为有限维、无边界、Hausdorff、第二可数的光滑流形，Lie 群也采用\chapref{g1:ch:16}的约定。

\subsection{主作用、局部截面与过渡函数}
\label{g1:subsec:1801:01}

在主丛的一条纤维中，选定一个参考点后，每个点都由唯一的群元素得到；参考点随底点光滑变化时，便形成局部乘积。局部截面负责选取这些参考点，过渡函数记录两套选择之间的群元素。

\begin{definition}[主丛]\label{g1:def:principal-bundle}
	设 \(G\) 为有限维实 Lie 群，\(P,M\) 为无边界、Hausdorff、第二可数的光滑流形。一个光滑\term[zhucong]{主丛}[principal bundle]由光滑满射 \(\pi:P\to M\) 和光滑右作用 \(P\times G\to P\)、\((p,g)\mapsto pg\) 组成，满足 \(\pi(pg)=\pi(p)\)，且每个 \(x\in M\) 都有开邻域 \(U\) 及微分同胚
	\[
		\Phi_U:U\times G\longrightarrow\pi^{-1}(U),
		\qquad \pi(\Phi_U(x,a))=x,\qquad
		\Phi_U(x,ag)=\Phi_U(x,a)g.
	\]
	称 \(G\) 为结构群，\(P\) 为总空间，\(M\) 为底空间。覆盖恒等映射的主丛同构是保持右作用的微分同胚；覆盖一般光滑映射的主丛映射也要求保持右作用与投影。
\end{definition}

局部模型说明每条纤维 \(P_x\) 是一条自由的 \(G\) 轨道：给定 \(p,q\in P_x\)，存在唯一 \(g\in G\) 使 \(q=pg\)。因此，在一个纤维中选定 \(p\) 后，\(g\mapsto pg\) 才把它识别为群 \(G\)；没有这个选择时，不能在 \(P_x\) 中直接把两个点相乘。局部平凡化提供光滑截面 \(s_U(x)=\Phi_U(x,e)\)，并满足 \(\Phi_U(x,g)=s_U(x)g\)。反过来，任意光滑局部截面都给出这样的平凡化：在已有平凡化中写 \(s_U(x)=\Phi_U^0(x,h(x))\)，则相应映射为 \((x,g)\mapsto\Phi_U^0(x,h(x)g)\)，其逆由左乘 \(h(x)^{-1}\) 给出。

\begin{proposition}[主丛与自由固有作用]\label{g1:prop:principal-free-proper}
	设 \(G\) 为有限维实 Lie 群，所有流形均无边界、Hausdorff、第二可数。主 \(G\) 丛 \(\pi:P\to M\) 的右作用自由且固有，纤维恰为轨道，并且 \(M\) 与轨道商 \(P/G\) 自然微分同胚。反之，\(G\) 在 \(P\) 上的任意光滑自由固有右作用都使商投影 \(P\to P/G\) 成为主 \(G\) 丛。
\end{proposition}
\begin{proof}
	自由性与纤维等于轨道已由局部模型得到。为证明固有性，设 \(p_j\to p\)、\(p_jg_j\to q\)。投影连续给出 \(\pi(p)=\pi(q)=x\)。选取 \(x\) 附近的主丛平凡化，把充分大的 \(j\) 对应的两点写成 \((x_j,a_j)\)、\((x_j,b_j)\)，则 \(a_j\to a\)、\(b_j\to b\)，并且 \(b_j=a_jg_j\)。故 \(g_j=a_j^{-1}b_j\to a^{-1}b\)。\cref{g1:prop:proper-action-criteria}的序列判据适用于右作用，只须把它改写为左作用 \(g*p=pg^{-1}\)；群元素取逆不会改变子列收敛条件。因此作用固有。

	\cref{g1:thm:free-proper-quotient}赋予 \(P/G\) 光滑结构。诱导的集合双射 \(P/G\to M\) 在任一局部乘积中都是 \([s_U(x)g]\mapsto x\)，它及其逆均光滑，因而为微分同胚。

	反过来，对左作用 \(g*p=pg^{-1}\) 应用\cref{g1:lem:free-proper-local-product}，得到横截面 \(S\) 与 \(G\times S\to SG\)、\((g,s)\mapsto sg^{-1}\) 的微分同胚。交换两个因子并取逆，便得到右作用所需的 \(S\times G\to SG\)、\((s,g)\mapsto sg\)。商定理又把 \(S\) 微分同胚地识别为底空间的开集，所以这正是主丛平凡化。
\end{proof}

\begin{example}[齐性空间上的主丛]\label{g1:exm:homogeneous-principal-bundle}
	设 \(H\) 为有限维实 Lie 群 \(G\) 的闭子群。\cref{g1:thm:homogeneous-space}不仅给出 \(G/H\) 的光滑结构，还给出局部乘积 \((\exp W)\times H\to(\exp W)H\) 及其左平移。因此 \(G\to G/H\) 是主 \(H\) 丛，右作用为群内的右乘。这里结构群是 \(H\)，不是 \(G\)；底空间也不必是一个群。
\end{example}

现在取一族局部截面 \(s_i:U_i\to P\)。交叠处，每条纤维内的自由传递性给出唯一函数 \(g_{ij}:U_i\cap U_j\to G\)，使
\begin{equation}\label{g1:eq:principal-transition}
	s_j(x)=s_i(x)g_{ij}(x),\qquad
	g_{ii}=e,\qquad g_{ij}g_{jk}=g_{ik}.
\end{equation}
在 \(s_i\) 给出的平凡化中，\(g_{ij}\) 就是 \(s_j\) 的群坐标，所以光滑。三重交叠上，\(s_k=s_jg_{jk}=s_i g_{ij}g_{jk}\)，再由自由性得到余循环等式；特别地 \(g_{ji}=g_{ij}^{-1}\)。同一点 \(p\in P_x\) 若写成 \(s_i(x)a_i=s_j(x)a_j\)，则
\begin{equation}\label{g1:eq:principal-coordinate-transition}
	a_i=g_{ij}(x)a_j.
\end{equation}
截面从 \(i\) 换到 \(j\) 时右乘 \(g_{ij}\)；表示同一点的群坐标从 \(j\) 换到 \(i\) 时，按上式左乘它。

\ProseParagraphBreak

如果换用 \(s_i'=s_i h_i\)，其中 \(h_i:U_i\to G\) 光滑，则
\begin{equation}\label{g1:eq:principal-transition-change}
	g_{ij}'=h_i^{-1}g_{ij}h_j.
\end{equation}
这是在同一个主丛上改选局部截面，底点 \(x\) 保持不变。局部联络形式也随截面改变，其变换律将同时涉及这里的群元素及其导数。

\subsection{由余循环重建主丛}
\label{g1:subsec:1801:02}

反过来，能否只给出 \(g_{ij}\)，就得到一个总空间？\cref{g1:thm:bundle-cocycle-gluing}已经对向量丛证明过这种粘合。这里用群中的左乘代替矩阵作用；需要保留的不只是局部公式，还包括不同纤维能够分离、总空间具有可数坐标基这些整体条件。

\begin{theorem}[主丛的余循环构造]\label{g1:thm:principal-cocycle-gluing}
	设 \(M\) 为有限维、无边界、Hausdorff、第二可数光滑流形，\(G\) 为有限维实 Lie 群，\(\{U_i\}_{i\in I}\) 为 \(M\) 的开覆盖。设光滑映射 \(g_{ij}:U_i\cap U_j\to G\) 满足 \(g_{ii}=e\)、\(g_{ij}g_{jk}=g_{ik}\)。则存在主 \(G\) 丛 \(\pi:P\to M\) 及局部截面 \(s_i\)，满足 \(s_j=s_i g_{ij}\)。它在保持这些局部截面的主丛同构意义下唯一。
\end{theorem}
\begin{proof}
	取不交并 \(\mathcal P=\coprod_i(U_i\times G)\)，规定
	\[
		(i,x,a)\sim(j,y,b)
		\quad\Longleftrightarrow\quad
		x=y\quad\text{且}\quad a=g_{ij}(x)b.
	\]
	自反性来自 \(g_{ii}=e\)，对称性来自 \(g_{ji}=g_{ij}^{-1}\)，传递性来自 \(g_{ij}g_{jk}=g_{ik}\)。令 \(P=\mathcal P/\!\sim\)，赋予商拓扑，定义 \(\pi([i,x,a])=x\)。关系没有改变 \(x\)，故 \(\pi\) 良定义且连续。

	商映射 \(Q:\mathcal P\to P\) 为开映射。若 \(O\subset\mathcal P\) 开，其饱和集在第 \(j\) 个副本中的部分为
	\[
		\bigcup_i\{(x,b):x\in U_i\cap U_j,\ (i,x,g_{ij}(x)b)\in O\}.
	\]
	每项是光滑同胚 \((x,b)\mapsto(x,g_{ij}(x)b)\) 下的开原像，所以饱和集开。于是
	\[
		\Phi_i:U_i\times G\longrightarrow\pi^{-1}(U_i),\qquad
		\Phi_i(x,a)=[i,x,a]
	\]
	是连续开双射，因而为同胚。双射性来自每个相应等价类可以唯一换算到第 \(i\) 个副本。它们的转换为
	\[
		\Phi_i^{-1}\Phi_j(x,b)=(x,g_{ij}(x)b),
	\]
	光滑且具有光滑逆。把 \(U_i\) 和 \(G\) 再限制到坐标域，便得到维数 \(\dim M+\dim G\) 的相容图册。

	若两点投影不同，用 \(M\) 中不交开邻域的逆像分离；若投影相同，把它们放进同一个开集 \(\pi^{-1}(U_i)\)，利用 \(U_i\times G\) 的 Hausdorff 性分离。故 \(P\) 为 Hausdorff 空间。第二可数性也要单独说明：从 \(M\) 的一个可数基中，对每个包含于某个 \(U_i\) 的基元素选择这样的 \(U_i\)，所选 \(U_i\) 构成可数子覆盖。各 \(U_i\times G\) 都有可数基，其像的可数并就给出 \(P\) 的可数基。因而任意开覆盖参与粘合也不会破坏总空间的第二可数性。

	定义 \([i,x,a]g=[i,x,ag]\)。若 \(a=g_{ij}b\)，则 \(ag=g_{ij}(bg)\)，故右作用良定义；在每个 \(\Phi_i\) 中就是群的右乘，因而光滑。\(\pi\) 局部为乘积投影，\(\Phi_i\) 为等变平凡化，故得到主丛。取 \(s_i(x)=[i,x,e]\)，则
	\[
		s_j(x)=[j,x,e]=[i,x,g_{ij}(x)]=s_i(x)g_{ij}(x).
	\]
	最后，若另一个主丛有同样关系的截面 \(t_i\)，公式 \([i,x,a]\mapsto t_i(x)a\) 良定义，在各局部乘积中是恒等坐标映射，故为主丛同构。保持指定截面与右作用的任何同构都必须使用这一公式，因此唯一。
\end{proof}

余循环本身仍依赖截面。如果两组余循环满足\eqref{g1:eq:principal-transition-change}，把旧余循环构造出的点 \([i,x,a]\) 送到新余循环构造中的 \([i,x,h_i(x)^{-1}a]\)，便得到主丛同构：由 \(a=g_{ij}b\) 可算出 \(h_i^{-1}a=g_{ij}'h_j^{-1}b\)，所以公式与代表无关，且局部光滑可逆。覆盖不同时，先取两组开集的交集作为共同细化即可比较。这里保持局部截面的唯一性，与更换截面后允许的主丛同构，是两种不同的陈述。

\begin{example}[有全局截面的主丛]\label{g1:exm:principal-global-section}
	主 \(G\) 丛 \(P\to M\) 与 \(M\times G\) 同构，当且仅当它有一个光滑全局截面 \(s\)。一方面，乘积中的 \(x\mapsto(x,e)\) 是截面；另一方面，\((x,g)\mapsto s(x)g\) 在每条纤维上双射，其在任何局部平凡化中的表达为 \((x,g)\mapsto(x,h(x)g)\)，故它与逆都光滑。局部截面在定义中总存在，全局截面的存在则是额外结论。向量丛即使非平凡也有零截面，主丛纤维没有相应的零点，因此不能用零截面的直觉推断这个条件。
\end{example}

\subsection{关联向量丛和标架丛}
\label{g1:subsec:1801:03}

主丛的一点可以理解为一种局部参照。若一个空间 \(F\) 上也规定了 \(G\) 怎样改变参照，就应能把 \(F\) 随着主丛一起粘到各个底点上。关键是：主丛中的参照乘上 \(g\) 后，\(F\) 中的分量必须按 \(g^{-1}\) 改变，才能表示同一个对象。

\begin{proposition}[关联丛]\label{g1:prop:associated-fiber-bundle}
	设 \(\pi:P\to M\) 为有限维实 Lie 群 \(G\) 的光滑主丛，\(F\) 为具有光滑左 \(G\) 作用的有限维光滑流形；所有流形均无边界、Hausdorff、第二可数。规定
	\[
		(p,y)\cdot g=(pg,g^{-1}\cdot y),\qquad
		P\times_G F\coloneqq(P\times F)/G.
	\]
	该商具有自然的光滑流形结构，投影 \([p,y]\mapsto\pi(p)\) 局部平凡，纤维为 \(F\)。若主丛截面满足 \(s_j=s_i g_{ij}\)，则局部表示 \((x,y_i)\mapsto[s_i(x),y_i]\) 的转换为 \(y_i=g_{ij}(x)\cdot y_j\)。称此丛为 \(P\) 的\term[guanliancong]{关联丛}[associated bundle]。
\end{proposition}
\begin{proof}
	先核对所写公式确为右作用：连续作用 \(g\) 再作用 \(h\)，第二因子变为 \(h^{-1}\cdot(g^{-1}\cdot y)=(gh)^{-1}\cdot y\)。它自由，因为第一因子的作用自由。它也固有：若 \((p_j,y_j)\to(p,y)\)，且 \((p_jg_j,g_j^{-1}\cdot y_j)\) 收敛，则主丛作用的固有性给出 \(g_j\) 的收敛子列，正好满足\cref{g1:prop:proper-action-criteria}的判据。于是\cref{g1:thm:free-proper-quotient}使商为无边界、Hausdorff、第二可数光滑流形，商映射为淹没。

	在一个主丛平凡化中，每个 \(p\) 唯一写成 \(s_i(x)a\)。定义
	\[
		[p,y]\longmapsto(x,a\cdot y),\qquad
		(x,z)\longmapsto[s_i(x),z].
	\]
	若 \((p,y)\) 换成 \((pg,g^{-1}\cdot y)\)，则第一式的群坐标由 \(a\) 变成 \(ag\)，输出仍为 \((x,a\cdot y)\)，所以良定义。两式互逆；第二式显然光滑，第一式与商映射的复合光滑，故由\cref{g1:prop:quotient-smooth-universal-property}也光滑。它们给出所需局部平凡化。最后
	\[
		[s_j(x),y_j]=[s_i(x)g_{ij}(x),y_j]
		=[s_i(x),g_{ij}(x)\cdot y_j],
	\]
	得到转换公式。
\end{proof}

\begin{theorem}[关联向量丛]\label{g1:thm:associated-vector-bundle}
	设 \(\pi:P\to M\) 为有限维实 Lie 群 \(G\) 的光滑主右丛，\(M,P\) 无边界、Hausdorff、第二可数，\(V\) 为有限维实或复向量空间，\(\rho:G\to\operatorname{GL}(V)\) 为光滑线性表示。取开覆盖 \(\{U_i\}\) 上的光滑局部截面 \(s_i:U_i\to P\)，并以 \(s_j=s_i g_{ij}\) 定义过渡函数 \(g_{ij}:U_i\cap U_j\to G\)。则
	\[
		E=P\times_\rho V\coloneqq(P\times V)/G,
		\qquad(p,v)\cdot g=(pg,\rho(g)^{-1}v)
	\]
	具有自然的秩 \(\dim_{\mathbb K}V\) 的光滑 \(\mathbb K\) 向量丛结构，其中 \(\mathbb K=\mathbb R\) 或 \(\mathbb C\)。局部坐标 \([s_i(x),v_i]\) 满足
	\begin{equation}\label{g1:eq:associated-vector-transition}
		v_i=\rho(g_{ij})v_j.
	\end{equation}
	对固定 \(p\in P_x\)，映射 \(v\mapsto[p,v]\) 是 \(V\to E_x\) 的线性同构。
\end{theorem}
\begin{proof}
	应用\cref{g1:prop:associated-fiber-bundle}，所得局部转换为 \((x,v_j)\mapsto(x,\rho(g_{ij}(x))v_j)\)。因为 \(\rho\) 是同态，这些可逆线性转换满足同样的余循环等式。由此，局部乘积中的加法、数乘在交叠处相容，给出各纤维与参照无关的向量空间结构，局部平凡化线性且光滑。复情形只把所有纤维运算和转换理解为复线性的，底空间及总空间的光滑结构仍按实流形理解。固定 \(p\) 后，每个类都可唯一写成 \([p,v]\)，故得到所述同构。
\end{proof}

这类构造也常称“伴随向量丛”。本书用“关联向量丛”表示一般 \(\rho\) 的情形，把伴随表示 \(\operatorname{Ad}\) 给出的特例记为 \(\operatorname{ad}(P)\)。表示不必忠实；关联构造可能丢掉主丛的一部分信息，不能从某一个关联向量丛自动恢复原主丛。

\begin{proposition}[截面与等变函数]\label{g1:prop:associated-sections}
	设 \(P\to M\) 为有限维实 Lie 群 \(G\) 的光滑主丛，\(\rho:G\to\operatorname{GL}(V)\) 为有限维实或复光滑表示，所有流形无边界、Hausdorff、第二可数。关联向量丛 \(E=P\times_\rho V\) 的光滑截面，与满足
	\begin{equation}\label{g1:eq:associated-equivariant-function}
		f(pg)=\rho(g^{-1})f(p)
	\end{equation}
	的光滑函数 \(f:P\to V\) 一一对应，对应式为 \(\sigma(x)=[p,f(p)]\)、\(p\in P_x\)。取开覆盖 \(\{U_i\}\) 上的光滑主丛截面 \(s_i\)，并以 \(s_j=s_i g_{ij}\) 定义过渡函数；其局部分量为 \(f_i=f\circ s_i\)，并满足 \(f_i=\rho(g_{ij})f_j\)。
\end{proposition}
\begin{proof}
	给定 \(\sigma\)，固定 \(p\in P_x\) 后，上一结果给出的线性同构使 \(\sigma(x)=[p,f(p)]\) 唯一确定 \(f(p)\)。换用 \(pg\) 时，\([pg,f(pg)]=[p,\rho(g)f(pg)]\)，故等于原截面值当且仅当\eqref{g1:eq:associated-equivariant-function}成立。在 \(p=s_i(x)a\) 的局部坐标中，\(f(p)=\rho(a^{-1})f_i(x)\)，因此截面光滑蕴含 \(f\) 光滑。

	反过来，等变关系使 \([p,f(p)]\) 与 \(p\in P_x\) 的选择无关，所以给出一个截面。它在 \(s_i\) 中的分量是光滑函数 \(f_i\)，因而光滑。两种构造互逆；而 \(s_j=s_i g_{ij}\) 给出 \(f_j=\rho(g_{ij}^{-1})f_i\)，等价于所写分量转换。
\end{proof}

\begin{proposition}[标架丛与向量丛]\label{g1:prop:frame-associated-equivalence}
	设 \(M\) 为有限维、无边界、Hausdorff、第二可数光滑流形，\(E\to M\) 为秩 \(r\geq1\) 的实光滑向量丛。\(\operatorname{Fr}(E)\to M\) 是主 \(\operatorname{GL}(r,\mathbb R)\) 丛，且标准表示给出自然向量丛同构
	\[
		\operatorname{Fr}(E)\times_{\operatorname{GL}(r,\mathbb R)}\mathbb R^r
		\longrightarrow E,\qquad [u,v]\longmapsto u(v).
	\]
	反之，若 \(P\to M\) 是主 \(\operatorname{GL}(r,\mathbb R)\) 丛，令 \(E=P\times_{\operatorname{GL}(r,\mathbb R)}\mathbb R^r\)，则
	\[
		P\longrightarrow\operatorname{Fr}(E),\qquad
		p\longmapsto u_p,\qquad u_p(v)=[p,v]
	\]
	是自然的主丛同构。复向量丛与 \(\operatorname{GL}(r,\mathbb C)\) 满足相同结论。
\end{proposition}
\begin{proof}
	\cref{g1:def:frame-bundle}已给出 \(\operatorname{Fr}(E)\) 的局部坐标 \(u=e_i(x)A\)，右作用在其中是 \((x,A)B=(x,AB)\)，所以满足主丛定义。第一式在等价代表下不变，因为 \((uA)(A^{-1}v)=u(v)\)。在参考标架 \(e_i\) 下，它把 \([e_i(x),v]\) 送到 \(e_i(x)v\)，局部就是恒等向量坐标，故光滑可逆并且逐纤维线性。

	对于第二式，\(u_p\) 是线性同构，且
	\[
		u_{pA}(v)=[pA,v]=[p,Av]=(u_p\circ A)(v),
	\]
	所以映射保持右作用。固定 \(p\in P_x\)，任何标架 \(u:\mathbb R^r\to E_x\) 都唯一写成 \(u_p\circ A\)，从而此映射在每条纤维上双射。在局部截面 \(s_i\) 给出的主丛坐标及标架 \(u_{s_i}\) 给出的标架丛坐标中，它为 \((x,A)\mapsto(x,A)\)，故为微分同胚。复情形逐步使用复线性同构，证明相同。
\end{proof}

\begin{definition}[伴随丛]\label{g1:def:adjoint-bundle}
	设 \(P\to M\) 为有限维实 Lie 群 \(G\) 的光滑主丛，\(\mathfrak g\) 为 \(G\) 的 Lie 代数。由伴随表示构造的实向量丛
	\[
		\operatorname{ad}(P)\coloneqq P\times_{\operatorname{Ad}}\mathfrak g
	\]
	称为\term[banshuicong]{伴随丛}[adjoint bundle]。其等价关系为 \((p,X)\sim(pg,\operatorname{Ad}_{g^{-1}}X)\)，取光滑局部截面 \(s_i:U_i\to P\)，以 \(s_j=s_i g_{ij}\) 定义过渡函数，并将局部元素写为 \([s_i(x),X_i(x)]\)。则局部分量满足 \(X_i=\operatorname{Ad}_{g_{ij}}X_j\)。
\end{definition}

每条纤维还带有 Lie 括号：用同一个 \(p\in P_x\) 表示两个元素，定义 \(\bigl[[p,X],[p,Y]\bigr]=[p,[X,Y]]\)，这里左侧外层括号表示纤维中的 Lie 括号。由于 \(\operatorname{Ad}_g[X,Y]=[\operatorname{Ad}_gX,\operatorname{Ad}_gY]\)，改变 \(p\) 后所得类不变；在局部坐标中就是 \(\mathfrak g\) 的固定双线性括号，故随底点光滑。若 \(G\) 交换，则所有 \(\operatorname{Ad}_g\) 都是恒等映射，于是 \([p,X]\mapsto(\pi(p),X)\) 给出典范同构 \(\operatorname{ad}(P)\cong M\times\mathfrak g\)，即使 \(P\) 本身可能并不平凡。后文的曲率会取值于这个伴随丛。它的纤维是带括号的向量空间 \(\mathfrak g\)，而原主丛的纤维是一条自由传递的群轨道。

\subsection{结构群约化及截面判据}
\label{g1:subsec:1801:04}

给向量丛指定内积后，正交单位标架之间只需正交矩阵换算。\cref{g1:def:frame-structure-reduction}已把这称为结构群约化。对于一般主 \(G\) 丛，问题变为：能否在每条 \(G\) 纤维中，光滑地选出一条固定闭子群 \(H\subseteq G\) 的轨道？选择单个点会使主丛平凡，而选择一条 \(H\) 轨道只要求减少允许的参照。

\begin{definition}[主丛的结构群约化]\label{g1:def:principal-reduction}
	设 \(P\to M\) 为有限维实 Lie 群 \(G\) 的光滑主丛，\(H\subseteq G\) 为闭子群，所有流形无边界、Hausdorff、第二可数。一个 \(H\) 结构群约化是 \(P\) 的光滑嵌入子流形 \(Q\)，它在右 \(H\) 作用下不变，且限制投影 \(Q\to M\) 与该作用组成主 \(H\) 丛。
\end{definition}

这里要求 \(H\) 闭，是为了使用\cref{g1:thm:homogeneous-space}中 \(G/H\) 的 Hausdorff 光滑结构及局部截面。我们将 \(P/H\) 定义为主丛总空间对右 \(H\) 作用的轨道商；它的一个点是一条 \(H\) 轨道，并非整条 \(G\) 纤维。仍有投影 \(r:P/H\to M\)、\(pH\mapsto\pi(p)\)。

\begin{theorem}[约化的截面判据]\label{g1:thm:reduction-section-criterion}
	设 \(\pi:P\to M\) 为有限维实 Lie 群 \(G\) 的光滑主丛，\(H\subseteq G\) 为闭子群，所有流形无边界、Hausdorff、第二可数。令 \(P/H\) 为 \(P\) 对右 \(H\) 作用的轨道商，赋予自然商流形结构，并定义 \(r:P/H\to M\)、\(r(pH)=\pi(p)\)。则 \(r\) 是纤维为 \(G/H\) 的光滑局部平凡丛，并有自然同构
	\[
		P\times_G(G/H)\longrightarrow P/H,\qquad
		[p,gH]\longmapsto pgH.
	\]
	主丛 \(P\) 的 \(H\) 约化与 \(r\) 的光滑截面一一对应，具体为
	\[
		Q\longmapsto\sigma_Q(x)=Q_x\in P/H,
		\qquad
		\sigma\longmapsto Q_\sigma=\{p\in P:pH=\sigma(\pi(p))\}.
	\]
	若 \(Q\) 为约化，则由 \(H\) 在 \(G\) 上的左乘构造的丛 \(Q\times_HG\) 具有自然主 \(G\) 丛结构，并有主丛同构
	\begin{equation}\label{g1:eq:extension-structure-group}
		Q\times_HG\longrightarrow P,\qquad[q,g]\longmapsto qg.
	\end{equation}
\end{theorem}
\begin{proof}
	先构造 \(P/H\)。右 \(H\) 作用自由。若 \(p_j\to p\)、\(p_jh_j\to q\)，主 \(G\) 作用的固有性给出 \(h_j\) 在 \(G\) 中的收敛子列；因为 \(H\) 闭，极限仍在 \(H\)，故 \(H\) 作用固有。商定理给出 \(P/H\) 的光滑结构。在主丛局部坐标中，商映射为
	\[
		U_i\times G\longrightarrow U_i\times(G/H),\qquad
		(x,g)\longmapsto(x,gH).
	\]
	\cref{g1:thm:homogeneous-space}保证 \(G\to G/H\) 有光滑局部截面，故上述商坐标和其逆光滑。交叠时 \(g_jH\) 变成 \(g_{ij}(x)g_jH\)，于是 \(r\) 局部平凡。公式 \([p,gH]\mapsto pgH\) 不随代表变化，因为 \((pa)(a^{-1}g)H=pgH\)；它在刚才的局部坐标中为恒等映射，所以给出所述自然同构。

	若已有 \(H\) 约化 \(Q\)，每个 \(Q_x\) 是唯一的一条 \(H\) 轨道，故 \(\sigma_Q\) 良定义。在 \(Q\) 的局部截面 \(t:U\to Q\) 中，\(\sigma_Q(x)=t(x)H\)，是光滑映射，且 \(r\sigma_Q=\operatorname{id}_M\)。

	反过来，给定光滑截面 \(\sigma\)，在 \(P\) 的局部截面 \(s_i\) 下，把它写成 \(\sigma(x)=(x,k(x)H)\)。这里首先确定的是一个光滑映射 \(U_i\to G/H\)，不能预先假定已有全局的 \(G\) 值函数 \(k\)。在任一点附近，使用 \(G\to G/H\) 的光滑局部截面，并适当缩小 \(U_i\)，才得到光滑提升 \(k:U_i\to G\)。于是
	\[
		Q_\sigma|_{U_i}=\{s_i(x)k(x)h:x\in U_i,\ h\in H\}.
	\]
	在 \(P|_{U_i}\cong U_i\times G\) 中，光滑坐标改变 \((x,g)\mapsto(x,k(x)^{-1}g)\) 把它送到 \(U_i\times H\)。闭子群 \(H\) 是嵌入 Lie 子群，故这些局部表达证明 \(Q_\sigma\) 是嵌入子流形，且 \((x,h)\mapsto s_i(x)k(x)h\) 给出主 \(H\) 丛的局部平凡化。它继承 \(P\) 的 Hausdorff 性与第二可数性。两种构造逐纤维互逆，故得到一一对应。

	最后，\(Q\times_HG\) 的关系为 \((q,g)\sim(qh,h^{-1}g)\)。右 \(G\) 作用定义为 \([q,g]a=[q,ga]\)，与等价关系相容；取 \(Q\) 的局部截面 \(t\)，则 \((x,g)\mapsto[t(x),g]\) 给出主 \(G\) 丛平凡化。\eqref{g1:eq:extension-structure-group}与代表无关，因为 \((qh)(h^{-1}g)=qg\)，并保持右 \(G\) 作用。在上述平凡化及 \(t\) 作为 \(P\) 的局部截面所给出的平凡化中，它为恒等坐标映射，故为主丛同构。
\end{proof}

由此还能直接读出过渡函数判据：\(P\) 有一个 \(H\) 约化，当且仅当经过细化覆盖，可以选取局部截面使全部 \(g_{ij}\) 取值于 \(H\)。若已有 \(Q\)，取它的局部截面即可；反过来，\(Q|_{U_i}=s_i(U_i)H\) 因 \(H\) 值过渡而相容，并由局部 \(U_i\times H\) 成为约化。这恢复了\chapref{g1:ch:12}以标架定义的结构群约化。

\begin{example}[度量选择的是一类标架]\label{g1:exm:metric-as-principal-reduction}
	设 \(E\to M\) 为秩 \(r\geq1\) 的实光滑向量丛。取 \(G=\operatorname{GL}(r,\mathbb R)\)、\(H=\operatorname{O}(r)\)。给定丛度量 \(h\)，全部正交单位标架组成 \(Q\subset\operatorname{Fr}(E)\)。\cref{g1:prop:metric-orthogonal-reduction}的光滑 Gram--Schmidt 构造给出局部正交单位标架，所以它是主 \(H\) 约化。对应的截面 \(M\to\operatorname{Fr}(E)/\operatorname{O}(r)\) 在每点选择的是“所有对 \(h_x\) 正交的标架”这一类，而不是某一组标架。

	同一事实可直接逐纤维核验：一个标架 \(u:\mathbb R^r\to E_x\) 确定内积 \(h_u(v,w)=\langle u^{-1}v,u^{-1}w\rangle\)，把 \(u\) 右乘正交矩阵不会改变它；反之，给出相同内积的两个标架恰相差一个正交矩阵。于是该商纤维就是 \(E_x\) 上所有正定内积的空间。度量总存在，并不意味着可以在整个底空间上选出一组光滑正交标架；后者还要求这个主 \(\operatorname{O}(r)\) 丛本身平凡。
\end{example}
% 纲要条目排为编号小节标题，并进入目录与书签；说明段仍保留于纲要注释中。

% !TeX root = ../../Book1.tex
% 纲要标题：### 18.2 主联络的等价描述
% 纲要位置：几何学计划纲要.md，第 818 行。
% 正文状态：初稿；尚未正式审核。

% 纲要内容开始
% 
% * 水平分布和等变性
% * 联络形式及再现条件
% * 局部联络形式的粘合
% * 主联络和关联丛上的协变微分
% 
% 纲要内容结束

\section{主联络的等价描述}
\label{g1:sec:1802}

主丛的一条纤维记录同一点上所有允许的标架或参考状态。沿底空间移动时，我们希望把一个状态连续地带到邻近纤维，但局部平凡化提供的“保持群坐标不变”会随所选截面改变。联络要做的是在总空间中规定哪些速度算作纯粹的底空间运动，再把这个规定写成可以计算、可以换局部截面的形式。本节固定有限维实 Lie 群 \(G\) 及光滑主右 \(G\) 丛 \(\pi:P\to M\)；所有流形无边界、Hausdorff 且第二可数。

\subsection{水平分布和等变性}
\label{g1:subsec:1802:01}

先分清两种速度。在 \(p\in P\) 处，沿纤维运动的速度组成
\[
	V_pP\coloneqq\ker(\mathrm{d}\pi)_p.
\]
局部平凡化把 \(\pi\) 写成 \(U\times G\to U\) 的投影，故这些核构成秩为 \(\dim G\) 的光滑子丛 \(VP\subset \mathrm{T}P\)，称为\term[chuizhifenbu]{垂直分布}[vertical distribution]。仅知道底空间速度 \(v\in \mathrm{T}_xM\)，还不能确定总空间速度：若 \(w\) 投影到 \(v\)，加上任意垂直向量仍投影到 \(v\)。

\begin{definition}\label{g1:def:principal-horizontal-connection}
	设 \(\pi:P\to M\) 为光滑主右 \(G\) 丛。一个光滑子丛 \(H\subset \mathrm{T}P\)，若对任意 \(p\in P\)、\(g\in G\) 满足
	\[
		\mathrm{T}_pP=H_p\oplus V_pP,\qquad
		(\mathrm{d}R_g)_pH_p=H_{pg},
	\]
	就称为一个\term[zhulianluo]{主联络}[principal connection]，其中 \(R_g(p)=pg\)。\(H\) 称为它的\term[shuipingfenbu]{水平分布}[horizontal distribution]，\(H_p\) 中的向量称为水平向量。
\end{definition}

直和条件使 \((\mathrm{d}\pi)_p|_{H_p}:H_p\to \mathrm{T}_{\pi(p)}M\) 成为线性同构。因此每个 \(v\in \mathrm{T}_{\pi(p)}M\) 有唯一水平向量 \(v_p^H\) 投影到它；这个向量称为 \(v\) 在 \(p\) 处的水平提升。其对 \(p,v\) 的依赖光滑，因为在局部丛标架中只是在求一个可逆矩阵的逆。等变条件进一步给出
\[
	(\mathrm{d}R_g)_p v_p^H=v_{pg}^H.
\]
它要求先改变纤维中的参考状态再提升，与先提升再作同一个右平移得到相同结果。任意垂直补空间未必满足这一条件；主联络要求这组选择与群作用相容。

\ProseParagraphBreak

垂直空间本身有一个来自群的自然坐标。对 \(X\in\mathfrak g=\mathrm{T}_eG\)，沿用基本向量场的构造，但本章使用右作用，记
\begin{equation}\label{g1:eq:principal-fundamental-field}
	X_p^\#\coloneqq\left.\frac{\mathrm{d}}{\mathrm{d}t}\right|_{t=0}p\exp(tX).
\end{equation}
这里使用右作用，基本场的括号因而带正号；第17章的左作用约定则带负号。

\begin{proposition}[主作用的垂直方向]\label{g1:prop:principal-fundamental-fields}
	设 \(G\) 为有限维实 Lie 群，\(\pi:P\to M\) 为光滑主右 \(G\) 丛。记 \(\mathfrak g=\operatorname{Lie}(G)\)，对 \(X\in\mathfrak g\) 令 \(X_p^\#=\left.\frac{\mathrm{d}}{\mathrm{d}t}\right|_{t=0}p\exp(tX)\)，并记 \(V_pP=\ker(\mathrm{d}\pi)_p\)、\(R_g(p)=pg\)。则映射 \(X\mapsto X_p^\#\) 是 \(\mathfrak g\to V_pP\) 的线性同构，并随 \(p\) 光滑变化。其基本场满足
	\begin{equation}\label{g1:eq:principal-fundamental-identities}
		(\mathrm{d}R_g)_pX_p^\#=(\operatorname{Ad}_{g^{-1}}X)_{pg}^\#,
		\qquad [X^\#,Y^\#]=[X,Y]^\#.
	\end{equation}
\end{proposition}
\begin{proof}
	取局部截面 \(s:U\to P\) 及平凡化 \(\Phi(x,h)=s(x)h\)。在这些坐标中，\(X^\#\) 对应 \((0,X_h^L)\)，其中 \(X_h^L=(\mathrm{d}L_h)_eX\) 是第16章的左不变场。因此它逐点给出所有垂直方向，而且 \(X\mapsto X_h^L\) 是随 \(h\) 光滑变化的同构。左不变场的括号定义了 \(\mathfrak g\) 的括号，故本章右作用的正号基本场保持括号。

	由\cref{g1:prop:exponential-naturality}的共轭公式，
	\[
		p\exp(tX)g=pg\exp(t\operatorname{Ad}_{g^{-1}}X).
	\]
	在零点求导便得到右平移公式。若要与\cref{g1:prop:fundamental-fields}比较，可把当前右作用改写为左作用 \(g\cdot p=pg^{-1}\)；那个左作用的基本场是 \(-X^\#\)，代入其反同态公式也得到此处的正号。
\end{proof}

\begin{example}\label{g1:exm:horizontal-additive-bundle}
	取加法群 \(G=\mathbb R\)、主丛 \(P=\mathbb R_x\times\mathbb R_t\to\mathbb R_x\)，右作用为 \((x,t)a=(x,t+a)\)。垂直方向由 \(\partial_t\) 张成。给定光滑函数 \(a(x)\)，令
	\[
		H_{(x,t)}=\operatorname{span}\{\partial_x-a(x)\partial_t\}.
	\]
	它与垂直方向互补，且右平移不改变 \(a(x)\)、\(\partial_x\) 和 \(\partial_t\)，所以是主联络。底空间单位速度的水平提升带有纤维分量 \(-a(x)\partial_t\)。水平不是图面上预先指定的“横着走”，而是由联络规定如何同时调整纤维状态。若把系数改为依赖 \(t\) 的一般函数 \(a(x,t)\)，仍有垂直补空间，却通常不再右平移不变。
\end{example}

\subsection{联络形式及再现条件}
\label{g1:subsec:1802:02}

有了水平与垂直的分解，任意速度 \(w\in \mathrm{T}_pP\) 都能写成“水平部分加 \(X_p^\#\)”。把这个 \(X\) 记下来，就得到一个 \(\mathfrak g\) 值一形式。这里 \(\Omega^k(P;\mathfrak g)\) 表示取值于固定向量空间 \(\mathfrak g\) 的交替 \(k\) 形式，等价于\cref{g1:def:bundle-valued-form}中平凡丛 \(P\times\mathfrak g\) 值的形式。取一组 \(\mathfrak g\) 的基后可逐分量拉回、求外微分；换一组常数基只作线性组合，故所得运算与基无关。

\begin{definition}\label{g1:def:principal-connection-form}
	设 \(G\) 为有限维实 Lie 群，\(\pi:P\to M\) 为光滑主右 \(G\) 丛，\(\mathfrak g=\operatorname{Lie}(G)\)。对 \(X\in\mathfrak g\) 令 \(X_p^\#=\left.\frac{\mathrm{d}}{\mathrm{d}t}\right|_{t=0}p\exp(tX)\)，并令 \(R_g(p)=pg\)。若光滑一形式 \(\omega\in\Omega^1(P;\mathfrak g)\) 对每个 \(X\in\mathfrak g\)、\(g\in G\) 都满足
	\begin{equation}\label{g1:eq:principal-connection-axioms}
		\omega_p(X_p^\#)=X,\qquad
		R_g^*\omega=\operatorname{Ad}_{g^{-1}}\omega
	\end{equation}
	则称其为\term[zhulianluoxingshi]{主联络形式}[principal connection form]。第一式称为再现条件；第二式的意思是 \(\omega_{pg}((\mathrm{d}R_g)_pw)=\operatorname{Ad}_{g^{-1}}(\omega_p(w))\)。
\end{definition}

再现条件固定了联络在垂直方向上的值，等变条件说明这些垂直坐标在右移后如何变化。\(\omega\) 不要求右平移不变：基本场自身按 \(\operatorname{Ad}_{g^{-1}}\) 变换，记录它的形式也必须作同样变换。只有在伴随作用平凡的情形，第二式才成为通常的不变性。

\begin{theorem}[主联络的两种描述]\label{g1:thm:principal-connection-equivalence}
	设 \(G\) 为有限维实 Lie 群，\(\pi:P\to M\) 为光滑主右 \(G\) 丛。对 \(X\in\operatorname{Lie}(G)\) 记 \(X_p^\#=\left.\frac{\mathrm{d}}{\mathrm{d}t}\right|_{t=0}p\exp(tX)\)。则右等变垂直补空间意义的主联络与主联络形式一一对应。对应由
	\[
		H_p=\ker\omega_p,\qquad
		w=w^H+\bigl(\omega_p(w)\bigr)_p^\#
	\]
	给出，其中 \(w^H\) 是 \(w\) 的水平部分。
\end{theorem}
\begin{proof}
	从 \(H\) 出发，沿光滑直和 \(\mathrm{T}P=H\oplus VP\) 投影到 \(VP\)，再用\cref{g1:prop:principal-fundamental-fields}的垂直同构取逆，得到光滑丛映射 \(\omega:\mathrm{T}P\to\mathfrak g\)。它逐纤维线性，因此是一形式。对 \(X_p^\#\)，垂直投影不作改变，故 \(\omega_p(X_p^\#)=X\)。若 \(w=w^H+X_p^\#\)，右平移把第一项送到水平空间，把第二项送到 \((\operatorname{Ad}_{g^{-1}}X)_{pg}^\#\)。于是右平移后的垂直坐标是 \(\operatorname{Ad}_{g^{-1}}X\)，正是等变条件。

	反向给定 \(\omega\)。再现条件使 \(\omega_p:\mathrm{T}_pP\to\mathfrak g\) 处处满射，故其核为光滑的秩 \(\dim M\) 子丛；这也可在局部矩阵中选可逆子块直接求核的光滑基。任意 \(w\) 都有分解
	\[
		w=\bigl(w-(\omega_p(w))_p^\#\bigr)+(\omega_p(w))_p^\#.
	\]
	第一项被 \(\omega_p\) 消去，第二项垂直。若一个垂直向量 \(X_p^\#\) 也在核中，再现条件给出 \(X=0\)，所以这个和是直和。等变条件将核送入核，应用 \(R_{g^{-1}}\) 又得到反向包含，故核右平移不变。两种构造都通过同一个唯一的水平—垂直分解恢复原对象，因此互为逆。
\end{proof}

\cref{g1:exm:horizontal-additive-bundle}的形式就是 \(\omega=\mathrm{d}t+a(x)\,\mathrm{d}x\)。它在 \(\partial_t\) 上取值为 \(1\)，在 \(\partial_x-a(x)\partial_t\) 上取值为零，因而把两种描述直接联系起来。

\ProseParagraphBreak

计算一般主丛的局部公式前，要把群方向上的微分写准确。回顾\cref{g1:prop:lie-exponential-differential}之前构造的形式
\[
	\theta_g^L=(\mathrm{d}L_{g^{-1}})_g:\mathrm{T}_gG\longrightarrow\mathfrak g.
\]
它称为左\term[maurercartanxingshi]{Maurer--Cartan 形式}[Maurer--Cartan form]。\footnote{Maurer，Ludwig Maurer（1859--1927），德国数学家，研究连续变换群与早期代数群理论。}上标强调采用左平移识别切空间；矩阵群中可以简写为 \(\theta^L=g^{-1}\mathrm{d}g\)，而一般 Lie 群始终用平移微分解释。

\begin{proposition}[群乘法的微分形式]\label{g1:prop:maurer-cartan-product}
	设 \(G\) 为有限维实 Lie 群，\(\theta^L_g=(\mathrm{d}L_{g^{-1}})_g:\mathrm{T}_gG\to\operatorname{Lie}(G)\) 为左 Maurer--Cartan 形式，\(m(g,h)=gh\)。对 \(v\in \mathrm{T}_gG\)、\(w\in \mathrm{T}_hG\)，有
	\begin{equation}\label{g1:eq:maurer-cartan-product}
		\theta^L_{gh}\bigl(\mathrm{d}m_{(g,h)}(v,w)\bigr)
		=\operatorname{Ad}_{h^{-1}}\theta^L_g(v)+\theta^L_h(w).
	\end{equation}
	特别地，\(L_a^*\theta^L=\theta^L\)、\(R_a^*\theta^L=\operatorname{Ad}_{a^{-1}}\theta^L\)。若 \(\rho:G\to\operatorname{GL}(V)\) 是有限维实或复表示，\(\rho_*=(\mathrm{d}\rho)_e\)，则
	\begin{equation}\label{g1:eq:maurer-cartan-representation}
		\rho_*\theta^L=\rho^{-1}\mathrm{d}\rho.
	\end{equation}
	右侧是 \(\operatorname{End}(V)\) 值一形式，表示先求 \(\rho\) 的微分，再左乘矩阵 \(\rho(g)^{-1}\)。
\end{proposition}
\begin{proof}
	乘法的微分为 \(\mathrm{d}R_h(v)+\mathrm{d}L_g(w)\)。写 \(v=(\mathrm{d}L_g)_eX\)，第一项是曲线 \(g\exp(tX)h=gh\exp(t\operatorname{Ad}_{h^{-1}}X)\) 的零点速度，应用 \(\theta^L_{gh}\) 后得到 \(\operatorname{Ad}_{h^{-1}}X\)。写 \(w=(\mathrm{d}L_h)_eY\)，第二项是 \(gh\exp(tY)\) 的速度，得到 \(Y\)。线性相加便得乘法公式。分别固定第一个或第二个因子，即得左右平移公式。

	同态满足 \(\rho\circ L_g=L_{\rho(g)}\circ\rho\)。由\cref{g1:prop:lie-homomorphism-differential}中这一恒等式的微分，\(\mathrm{d}\rho_g((\mathrm{d}L_g)_eX)=\rho(g)\rho_*(X)\)。左乘 \(\rho(g)^{-1}\)，便得最后一式。
\end{proof}

沿任何一条主纤维取轨道参数 \(\alpha_p(g)=pg\)，再现条件给出 \(\alpha_p^*\omega=\theta^L\)：因为 \(\mathrm{d}\alpha_p(X_g^L)=X_{pg}^\#\)。因此 \(\omega\) 的纤维方向部分已经完全固定；可以自由选择的，是底空间方向怎样与这些垂直方向混合。

\subsection{局部联络形式的粘合}
\label{g1:subsec:1802:03}

取主丛的局部截面 \(s_i:U_i\to P\)，令 \(\Phi_i(x,g)=s_i(x)g\)。把联络形式限制到截面上，得到 \(\mathfrak g\) 值一形式 \(A_i=s_i^*\omega\)。它只在底空间 \(U_i\) 上，通常比定义在整个 \(P\) 上的 \(\omega\) 更便于计算。若 \(G\) 是矩阵群，\(A_i\) 是矩阵值一形式，因此也称局部联络矩阵；一般 \(G\) 的 \(A_i\) 仍是 \(\mathfrak g\) 值形式，不必先选择一个矩阵实现。

\begin{theorem}[主联络的局部数据]\label{g1:thm:principal-connection-local-data}
	设 \(G\) 为有限维实 Lie 群，\(\pi:P\to M\) 为光滑主右 \(G\) 丛，\(\mathfrak g=\operatorname{Lie}(G)\)。取开覆盖 \(\{U_i\}\) 上的光滑局部截面 \(s_i:U_i\to P\)，以 \(s_j=s_i g_{ij}\) 定义过渡函数。令 \(\Phi_i:U_i\times G\to P|_{U_i}\)、\(\Phi_i(x,g)=s_i(x)g\)，并令 \(\theta^L_g=(\mathrm{d}L_{g^{-1}})_g\) 为左 Maurer--Cartan 形式。任意主联络形式 \(\omega\) 由 \(A_i=s_i^*\omega\) 唯一确定，且
	\begin{equation}\label{g1:eq:principal-connection-trivialization}
		(\Phi_i^*\omega)_{(x,g)}(v,\xi)
		=\operatorname{Ad}_{g^{-1}}\bigl((A_i)_x(v)\bigr)+\theta_g^L(\xi).
	\end{equation}
	在交叠域上，这些局部形式满足
	\begin{equation}\label{g1:eq:principal-connection-change}
		A_j=\operatorname{Ad}_{g_{ij}^{-1}}A_i+g_{ij}^*\theta^L.
	\end{equation}
	反过来，满足此规律的一族 \(A_i\in\Omega^1(U_i;\mathfrak g)\) 通过\eqref{g1:eq:principal-connection-trivialization}唯一地粘成主联络形式。
\end{theorem}
\begin{proof}
	\(\mathrm{d}\Phi_i(v,0)\) 是 \((\mathrm{d}s_i)v\) 经 \(R_g\) 推前的向量，等变条件使其 \(\omega\) 值等于 \(\operatorname{Ad}_{g^{-1}}A_i(v)\)。写 \(\xi=(\mathrm{d}L_g)_eX\)，则 \(\mathrm{d}\Phi_i(0,\xi)=X_{s_i(x)g}^\#\)，再现条件给出 \(X=\theta_g^L(\xi)\)。相加得到局部表示，因此它必然唯一。

	截面 \(s_j\) 在第 \(i\) 个平凡化中的表达为 \(x\mapsto(x,g_{ij}(x))\)。把这个映射代入局部表示，便得到\eqref{g1:eq:principal-connection-change}。

	反向给定 \(A_i\)，先在每个 \(U_i\times G\) 上定义右端的一形式。基本场为 \((0,X_g^L)\)，所以再现条件成立。右乘固定的 \(h\) 后，第一项变为 \(\operatorname{Ad}_{h^{-1}}\operatorname{Ad}_{g^{-1}}A_i\)，第二项由\cref{g1:prop:maurer-cartan-product}也变为 \(\operatorname{Ad}_{h^{-1}}\theta^L\)，因而等变条件成立。

	交叠处的坐标变换是 \(\Psi_{ij}(x,h)=(x,g_{ij}(x)h)\)，满足 \(\Phi_j=\Phi_i\circ\Psi_{ij}\)。记 \(\widetilde\omega_i=\operatorname{Ad}_{g^{-1}}A_i+\theta^L\)，乘法公式给出
	\begin{align*}
		\Psi_{ij}^*\widetilde\omega_i
		&=\operatorname{Ad}_{h^{-1}}\bigl(\operatorname{Ad}_{g_{ij}^{-1}}A_i+g_{ij}^*\theta^L\bigr)+\theta^L\\
		&=\operatorname{Ad}_{h^{-1}}A_j+\theta^L.
	\end{align*}
	这里各项均已拉回到交叠域与 \(G\) 的乘积，省略了底空间投影的记号。结果正是第 \(j\) 张图的形式，所以它们在总空间重叠部分相等，粘合为光滑 \(\omega\)。两条联络条件是局部可验证的，故粘合后仍成立。
\end{proof}

在矩阵群中，以上两式分别是
\[
	\Phi_i^*\omega=g^{-1}A_i g+g^{-1}\mathrm{d}g,
	\qquad A_j=g_{ij}^{-1}A_i g_{ij}+g_{ij}^{-1}\mathrm{d}g_{ij}.
\]
第二式恰是\cref{g1:prop:connection-local-data}中的换标架公式。现在可以看见其中两项的不同来源：共轭来自纤维参考状态的改变，附加微分项来自新截面本身随基点移动。局部截面 \(s_i\) 的全部切向量水平当且仅当 \(A_i=0\)。这样的截面比局部平凡化多满足一个微分条件；下一节将由曲率判断它何时存在。

\begin{example}\label{g1:exm:principal-circle-local-change}
	在平凡主圆周丛 \(U\times\operatorname{U}(1)\) 上，\(\mathfrak u(1)=i\mathbb R\)。选一个实一形式 \(a\)，令 \(A=i a\)，则
	\[
		\omega=i a+z^{-1}\mathrm{d}z.
	\]
	因为圆周群交换，伴随项没有共轭变化。若换截面 \(s'=s e^{i\varphi(x)}\)，其中 \(\varphi\) 为光滑实函数，则 \(A'=i(a+\mathrm{d}\varphi)\)。这只是同一个 \(\omega\) 的新表达。一个一般的 \(U\to\operatorname{U}(1)\) 映射未必有全局实值角度 \(\varphi\)，但 \(g^{-1}\mathrm{d}g\) 总有定义，所以用群值过渡函数写的公式比角度公式具有更大的适用范围。
\end{example}

变换律中的附加微分项说明，各个截面下的零形式一般不相容。整体存在性可以沿用单位分解的方法；这次把局部联络形式拉到总空间后加权，权重之和为一使其保留竖直方向的再现条件。

\begin{theorem}[主联络的存在]\label{g1:thm:principal-connection-existence}
	设 \(G\) 为有限维实 Lie 群，\(M\) 为无边界、Hausdorff、第二可数光滑流形。每个光滑主右 \(G\) 丛 \(P\to M\) 都有光滑主联络，无需假定 \(G\) 或 \(M\) 紧。
\end{theorem}
\begin{proof}
	取局部截面覆盖 \(\{U_i\}\)，在 \(\pi^{-1}(U_i)\) 上令 \(A_i=0\)，用\eqref{g1:eq:principal-connection-trivialization}得到局部主联络 \(\omega_i\)。这些 \(\omega_i\) 不要求在交叠处相同。由\cref{g1:thm:smooth-partition-unity}取从属单位分解 \(\{\chi_i\}\)，支集包含于 \(U_i\)、支集族局部有限、总和为 \(1\)。定义
	\[
		\omega=\sum_i(\chi_i\circ\pi)\omega_i,
	\]
	每项在 \(\pi^{-1}(U_i)\) 外补零。因为 \(\operatorname{supp}\chi_i\subset U_i\)，补零后在该开集外每点附近恒为零，故光滑；底空间支集族局部有限，拉回后的求和也局部有限。

	每个 \(\chi_i\circ\pi\) 右平移不变，故等变条件可逐项相加。再现条件为
	\[
		\omega(X^\#)=\sum_i(\chi_i\circ\pi)X=X.
	\]
	因此 \(\omega\) 是主联络形式。这里没有对非紧群作平均；起作用的是底空间的单位分解以及联络公理对权重和为 \(1\) 的加法稳定。
\end{proof}

\begin{proposition}[主联络的仿射空间]\label{g1:prop:principal-connection-affine}
	设 \(P\to M\) 为上述光滑主右 \(G\) 丛。任意两个主联络形式之差 \(\beta=\omega_1-\omega_0\) 在垂直方向为零，并满足 \(R_g^*\beta=\operatorname{Ad}_{g^{-1}}\beta\)。这些形式与 \(\Omega^1(M;\operatorname{ad}P)\) 自然一一对应。因此主联络的集合是以 \(\Omega^1(M;\operatorname{ad}P)\) 为方向空间的非空仿射空间。
\end{proposition}
\begin{proof}
	两个再现条件相减给出 \(\beta(X^\#)=0\)，两个等变条件相减给出所述变换式。回顾\cref{g1:def:adjoint-bundle}，伴随丛为 \(\operatorname{ad}P=P\times_{\operatorname{Ad}}\mathfrak g\)，满足 \([pg,Y]=[p,\operatorname{Ad}_gY]\)。对 \(v\in \mathrm{T}_xM\)，选 \(p\in P_x\) 及任意 \(w\in \mathrm{T}_pP\) 使 \(\mathrm{d}\pi(w)=v\)，令
	\[
		B_x(v)=[p,\beta_p(w)].
	\]
	同一点 \(p\) 处改变提升 \(w\) 只增加垂直向量，不影响值；把 \(p\) 改为 \(pg\)，可取 \(\mathrm{d}R_g(w)\)，此时
	\[
		[pg,\beta_{pg}(\mathrm{d}R_gw)]
		=[pg,\operatorname{Ad}_{g^{-1}}\beta_p(w)]
		=[p,\beta_p(w)].
	\]
	所以 \(B\) 与选择无关。在局部截面中它由 \(s_i^*\beta\) 表示，因而光滑且逐点对 \(v\) 线性。

	反向给定 \(B\in\Omega^1(M;\operatorname{ad}P)\)，对每个 \(p\in P_x\)、\(w\in \mathrm{T}_pP\)，由线性同构 \(Y\mapsto[p,Y]\) 唯一解出 \(\beta_p(w)\)，使 \([p,\beta_p(w)]=B_x(\mathrm{d}\pi(w))\)。局部表达保证光滑，定义直接保证消去垂直方向以及所需等变性。故 \(\omega_0+\beta\) 保留两条联络公理，也是主联络。两种对应互为逆，且任何两个联络之间有唯一的差；存在定理使集合非空，于是得到所称的仿射空间。
\end{proof}

这与\cref{g1:prop:connection-affine-space}的向量丛结论相呼应：单个联络包含一个规定好的微分部分，两联络之差才是底空间上的丛值一形式。\(\omega\) 本身不能直接下降成 \(\operatorname{ad}P\) 值一形式，因为它在垂直方向上的再现值通常非零。

\subsection{主联络和关联丛上的协变微分}
\label{g1:subsec:1802:04}

设 \(\rho:G\to\operatorname{GL}(V)\) 为有限维实或复表示，\(E=P\times_\rho V\)。\cref{g1:prop:associated-sections}把 \(E\) 的截面 \(\sigma\) 与等变函数 \(f:P\to V\) 对应：
\[
	\sigma(x)=[p,f(p)],\qquad
	f(pg)=\rho(g)^{-1}f(p).
\]
普通微分 \(\mathrm{d}f\) 还会记录沿纤维改变参考状态所造成的分量变化。协变微分只沿水平提升求导，正好排除这部分变化。

\begin{theorem}[关联丛上的诱导联络]\label{g1:thm:associated-covariant-derivative}
	设 \(G\) 为有限维实 Lie 群，\(\pi:P\to M\) 为光滑主右 \(G\) 丛，\(\omega\) 为主联络，\(\rho:G\to\operatorname{GL}(V)\) 为有限维实或复光滑表示。令 \(E=P\times_\rho V\)，采用 \((p,v)\sim(pg,\rho(g)^{-1}v)\) 的商约定，并记 \(\rho_*=(\mathrm{d}\rho)_e\)。对 \(\sigma\in\Gamma(E)\)，令 \(f:P\to V\) 为满足 \(f(pg)=\rho(g)^{-1}f(p)\)、\(\sigma(x)=[p,f(p)]\) 的光滑函数。对 \(v\in\mathrm{T}_xM\)、\(p\in P_x\)，以 \(v_p^H\) 表示满足 \((\mathrm{d}\pi)_p(v_p^H)=v\)、\(\omega_p(v_p^H)=0\) 的水平提升。则公式
	\begin{equation}\label{g1:eq:associated-horizontal-derivative}
		(\nabla\sigma)_x(v)=[p,(\mathrm{d}f)_p(v_p^H)],
		\qquad p\in P_x,
	\end{equation}
	定义与 \(p\) 的选择无关的丛联络。在局部截面 \(s_i\) 下，令 \(f_i=f\circ s_i\)、\(A_i=s_i^*\omega\)，则
	\begin{equation}\label{g1:eq:associated-local-derivative}
		\nabla\sigma=[s_i,\mathrm{d}f_i+\rho_*(A_i)f_i].
	\end{equation}
	这里的方括号对每个切向输入求值；实情形得到实联络，复情形得到对截面复线性的复联络。
\end{theorem}
\begin{proof}
	水平提升满足 \(v_{pg}^H=\mathrm{d}R_g(v_p^H)\)，等变函数满足 \(f\circ R_g=\rho(g)^{-1}f\)。对后一式求导，得到
	\[
		(\mathrm{d}f)_{pg}(v_{pg}^H)=\rho(g)^{-1}(\mathrm{d}f)_p(v_p^H).
	\]
	按关联丛的等价关系，两边与 \(pg,p\) 组合后表示同一纤维向量，证明选择无关。光滑性可在任意局部截面下检验。对底空间函数 \(h\)，截面 \(h\sigma\) 对应 \((h\circ\pi)f\)；在水平向量上使用普通乘积法则，得到 \(\nabla(h\sigma)=\mathrm{d}h\otimes\sigma+h\nabla\sigma\)。方向输入的函数线性及截面输入的实或复线性由微分直接继承，故确为联络。

	为求局部表达，先沿垂直基本场微分等变函数。指数自然性给出
	\[
		f(p\exp(tX))=\rho(\exp(tX))^{-1}f(p),
		\qquad (\mathrm{d}f)_p(X_p^\#)=-\rho_*(X)f(p).
	\]
	把 \(w=w^H+(\omega_p(w))_p^\#\) 代入，便得
	\begin{equation}\label{g1:eq:equivariant-covariant-derivative}
		(\mathrm{d}f)_p(w^H)=(\mathrm{d}f)_p(w)+\rho_*(\omega_p(w))f(p).
	\end{equation}
	因此总空间上的表达为 \(\mathrm{d}f+\rho_*(\omega)f\)，它消去垂直方向。拉回到 \(s_i\) 就得到\eqref{g1:eq:associated-local-derivative}。

	还可直接核对局部粘合。令 \(r_{ij}=\rho(g_{ij})\)。由\cref{g1:prop:lie-homomorphism-differential}的伴随交织公式及\eqref{g1:eq:maurer-cartan-representation}，
	\[
		\rho_*(A_j)=r_{ij}^{-1}\rho_*(A_i)r_{ij}+r_{ij}^{-1}\mathrm{d}r_{ij}.
	\]
	而 \(f_i=r_{ij}f_j\)，故乘积法则给出
	\[
		\mathrm{d}f_i+\rho_*(A_i)f_i
		=r_{ij}\bigl(\mathrm{d}f_j+\rho_*(A_j)f_j\bigr).
	\]
	这正是关联丛值一形式的分量变换，也说明局部公式与第15章的联络定义一致。
\end{proof}

\begin{corollary}[伴随丛的协变微分]\label{g1:cor:adjoint-connection-local}
	设 \(G\) 为有限维实 Lie 群，\(\pi:P\to M\) 为光滑主右 \(G\) 丛，\(\omega\) 为主联络，\(\mathfrak g=\operatorname{Lie}(G)\)，\(\operatorname{ad}P=P\times_{\operatorname{Ad}}\mathfrak g\)。取光滑局部截面 \(s_i:U_i\to P\)，令 \(A_i=s_i^*\omega\)，并把 \(\xi\in\Gamma(\operatorname{ad}P)\) 写为 \(\xi|_{U_i}=[s_i,\xi_i]\)。则 \(\omega\) 在 \(\operatorname{ad}P\) 上诱导联络，其局部表达为
	\[
		\nabla^{\operatorname{ad}}\xi=[s_i,\mathrm{d}\xi_i+[A_i,\xi_i]],
		\qquad [A_i,\xi_i](v)=[A_i(v),\xi_i].
	\]
\end{corollary}
\begin{proof}
	在上一定理中取 \(\rho=\operatorname{Ad}\)。\cref{g1:prop:adjoint-differential-bracket}证明了 \(\rho_*(X)=\operatorname{ad}_X\)，且 \(\operatorname{ad}_X(Y)=[X,Y]\)，代入即得公式。
\end{proof}

从主联络到关联丛联络的方向对所有表示都成立，反方向则取决于表示能看见多少信息。标架丛给出一个确切的双向对应。

\begin{proposition}[标架丛与线性联络]\label{g1:prop:frame-connection-equivalence}
	设 \(E\to M\) 为秩 \(r\geq1\) 的实光滑向量丛，其标架丛为 \(\operatorname{Fr}(E)\)，结构群为 \(\operatorname{GL}(r,\mathbb R)\)。\(\operatorname{Fr}(E)\) 上的主联络与 \(E\) 上的线性联络自然一一对应，正向对应为标准表示下的诱导联络。正有限秩复向量丛与 \(\operatorname{GL}(r,\mathbb C)\) 主标架丛有同样结论，其中线性联络取复线性。
\end{proposition}
\begin{proof}
	\cref{g1:prop:frame-associated-equivalence}给出 \(\operatorname{Fr}(E)\times_{\operatorname{GL}(r)}\mathbb F^r\cong E\)，映射为 \([u,v]\mapsto u(v)\)，其中 \(\mathbb F=\mathbb R\) 或 \(\mathbb C\)。因此主联络通过标准表示诱导 \(E\) 上的联络。

	反向，给定 \(E\) 上的联络 \(\nabla\)，在局部行标架 \(e_i\) 中取其联络矩阵 \(A_i\)。\(e_i\) 同时是 \(\operatorname{Fr}(E)\) 的局部截面，且 \(e_j=e_i g_{ij}\)。\cref{g1:prop:connection-local-data}已证明
	\[
		A_j=g_{ij}^{-1}A_i g_{ij}+g_{ij}^{-1}\mathrm{d}g_{ij}.
	\]
	复情形逐分量应用同一乘积法则，所得公式不变。\(\mathfrak{gl}(r,\mathbb F)\) 就是所有 \(r\times r\) 的 \(\mathbb F\) 矩阵，所以这些 \(A_i\) 均是合格的 Lie 代数值局部数据。由\cref{g1:thm:principal-connection-local-data}，它们唯一地构造出主联络。标准表示的微分是矩阵上的恒等映射，故再诱导回 \(E\) 后，局部矩阵仍为原来的 \(A_i\)，恢复 \(\nabla\)。反过来，从主联络得到 \(A_i\) 后再按局部公式重建，也由唯一性恢复原主联络。
\end{proof}

\begin{example}\label{g1:exm:adjoint-loses-central-connection}
	取任意主 \(\operatorname{U}(1)\) 丛。圆周群交换，所以 \(\operatorname{Ad}\) 是平凡表示，\(\operatorname{ad}P\) 自然等于 \(M\times i\mathbb R\)。无论原主联络是什么，诱导在伴随丛上的联络都是普通微分 \(\mathrm{d}\)，因为 \([A_i,\xi_i]=0\)。例如在 \(\mathbb R^2\times\operatorname{U}(1)\) 上，局部形式 \(A_0=0\) 与 \(A_1=i x\,\mathrm{d}y\) 确定不同主联络，却有完全相同的伴随丛联络。因此“主联络诱导伴随丛联络”并不意味着仅凭后者总能恢复前者。
\end{example}

一般表示中，两个主联络的差 \(B\in\Omega^1(M;\operatorname{ad}P)\) 在关联丛上变成 \(\rho_*(B)\)。若 \(\rho_*\) 单射，这些差仍可被区分；若 \(\rho_*\) 有核，就可能丢失方向。即使 \(\rho_*\) 单射，也不能说关联丛上的每个联络都来自原结构群：其局部矩阵还必须落在 \(\rho_*(\mathfrak g)\) 中。标架丛的标准表示之所以给出全体联络的双向对应，正是因为这个像等于所有矩阵。

% !TeX root = ../../Book1.tex
% 纲要标题：### 18.3 水平提升与曲率
% 纲要位置：几何学计划纲要.md，第 825 行。
% 正文状态：初稿；尚未正式审核。

% 纲要内容开始
% 
% * 路径提升及完备性范围
% * 曲率结构方程
% * Bianchi 恒等式的主丛形式
% * 小环路和平行移动的变化
% 
% 纲要内容结束

\section{水平提升与曲率}
\label{g1:sec:1803}

水平空间规定了怎样离开一个纤维而不额外作群方向的运动。沿底空间的一条路径不断采用这个方向，就应当得到两个端点纤维之间的比较。这里有两个需要分别回答的问题：所给路径能否完整提升，以及沿不同路径得到的比较为何会不同。前者由主作用的等变性控制，后者由曲率控制。

\ProseParagraphBreak

本节设 \(\pi:P\to M\) 为有限维实 Lie 群 \(G\) 的光滑右主丛，流形均有限维、Hausdorff、第二可数且无边界，\(\omega\in\Omega^1(P;\mathfrak g)\) 为主联络形式，\(H=\ker\omega\)。沿用上一节的约定：\(X^\#_p=\left.\frac{\mathrm{d}}{\mathrm{d}t}\right|_0p\exp(tX)\)、\(\omega(X^\#)=X\)，且 \(R_g^*\omega=\operatorname{Ad}_{g^{-1}}\omega\)。若 \(s:U\to P\) 为局部截面，则 \(A=s^*\omega\)，并以 \(p=s(x)g\) 表示局部纤维坐标。

\subsection{路径提升及完备性范围}
\label{g1:subsec:1803:01}

设 \(\gamma:[a,b]\to U\) 光滑，将它的一条提升写成 \(q(t)=s(\gamma(t))g(t)\)。由\eqref{g1:eq:principal-connection-trivialization}，水平条件正好是
\[
	0=\omega(q')=\operatorname{Ad}_{g^{-1}}A(\gamma')+\theta^L_g(g'),
\]
其中 \(\theta^L_g=(\mathrm{d}L_{g^{-1}})_g\)。把这个等式左平移回 \(\mathrm{T}_gG\)，得到
\begin{equation}\label{g1:eq:principal-horizontal-ode}
	g'(t)=-(\mathrm{d}R_{g(t)})_eB(t),\qquad B(t)=A_{\gamma(t)}(\gamma'(t)).
\end{equation}
矩阵群中它写成 \(g'=-Bg\)。右边是随时间变化的右不变向量场；“右不变”指对每个固定时间作右平移后不变，并不是说系数 \(B(t)\) 与时间无关。这个特征使我们能够在非紧 Lie 群上排除有限时间爆破。

\begin{lemma}[右不变方程的全区间解]\label{g1:lem:right-invariant-time-ode}
	设 \(G\) 为有限维实 Lie 群，\(\mathfrak g=\operatorname{Lie}(G)\)，\(B:[a,b]\to\mathfrak g\) 光滑，端点处按局部光滑延拓理解。对任意 \(t_0\in[a,b]\) 和 \(g_0\in G\)，方程
	\[
		g'(t)=-(\mathrm{d}R_{g(t)})_eB(t),\qquad g(t_0)=g_0
	\]
	（即\eqref{g1:eq:principal-horizontal-ode}）在整个 \([a,b]\) 上存在唯一解。解光滑依赖初始时刻、初值及 \(B\) 所含的有限维光滑参数。
\end{lemma}
\begin{proof}
	\cref{g1:thm:ode-local-existence-uniqueness}在单位元的一张坐标图中给出局部解。对每个 \(\tau\in[a,b]\)，该定理还给出初始时刻的一个邻域，以及从单位元出发都可使用的一段共同时间。用有限多个这样的邻域覆盖紧区间 \([a,b]\)，取各共同时间的一个更小正下界 \(\delta\)。于是从任意 \(\tau\in[a,b]\) 和单位元出发，都能至少求解到 \(|t-\tau|<\delta\) 与给定时间区间的交；\(\delta\) 不依赖 \(\tau\)。

	若 \(u(t)\) 是从单位元出发的这条短时间解，则 \(u(t)g_0\) 满足同一方程，因为
	\[
		\frac{\mathrm{d}}{\mathrm{d}t}(u g_0)
		=-(\mathrm{d}R_{g_0})_u(\mathrm{d}R_u)_eB
		=-(\mathrm{d}R_{u g_0})_eB.
	\]
	因此相同的时间 \(\delta\) 适用于任意 \(g_0\)，即使这些初值不在同一个紧集中。把 \([a,b]\) 分成长度小于 \(\delta/2\) 的有限小段，用前一段的终值作为后一段的初值，便可向前、向后走完整个区间。局部唯一性使重叠部分一致，也证明所得全区间解唯一。

	参数依赖由\cref{g1:thm:ode-smooth-dependence}及上述有限次拼接得到：固定一组参数后，每一步在一个共同局部解邻域内光滑依赖起始数据；缩小参数邻域，有限步的中间值仍落在相应邻域中。有限次复合保持光滑性。初始时刻和初值也包括在该定理的光滑变量内。
\end{proof}

每一步都能把初值右平移回同一个单位元附近，因此同一步长可以重复使用，直至走完整个紧参数区间。这个论证适用于一般的 \(G\) 与 \(M\)。

\begin{theorem}[主丛的水平提升]\label{g1:thm:principal-horizontal-lift}
	设 \(P\to M\) 为有限维实 Lie 群的光滑右主丛，\(M\) 无边界，\(\omega\) 为主联络。对任意连续、有限分片光滑路径 \(\gamma:[a,b]\to M\)、任意 \(t_0\in[a,b]\) 和 \(p_0\in P_{\gamma(t_0)}\)，存在唯一连续、分片光滑提升 \(q:[a,b]\to P\)，使下列等式在每个光滑片上成立：
	\[
		\pi\circ q=\gamma,\qquad q(t_0)=p_0,\qquad \omega(q')=0
	\]
	称这条提升为具有该初值的\term[shuipingtisheng]{水平提升}[horizontal lift]。若路径及初值光滑依赖有限维参数，并具有共同有限光滑分割，则提升在各片内都光滑依赖这些参数，端点运输在端点丛图中的表达也光滑依赖这些参数。
\end{theorem}
\begin{proof}
	先将光滑路径的紧像用有限个主丛平凡化域覆盖。时间区间的 Lebesgue 数性质给出有限分割，使每段闭区间的像包含在一个平凡化域内。在每段解\eqref{g1:eq:principal-horizontal-ode}，前一引理保证这一段可以完整走完。接续时将实际终点写入下一张丛图，作为新的初值。

	局部公式是内在条件 \(\omega(q')=0\) 的等价表达，因此换丛图不会改变所求对象。若两条提升具有同一初值，局部方程的唯一性先使它们在初始段相同，再依次使全部相邻段相同。这同时证明细分与丛图的选择无关。分片光滑路径只需把它的有限分段点加入上述分割；拐角处保留连续性，各片分别满足水平条件。

	对参数族，先在一个固定参数处选择有限分割，并让各段紧像留在所选丛图的内部。连续性与时间区间的紧性保证附近参数可以沿用同一批丛图和分割。局部方程的系数、各段末端的坐标变换与初值都光滑依赖参数，前一引理的参数结论及有限次复合遂给出所述光滑依赖。
\end{proof}

\begin{proposition}[主丛运输的运算规律]\label{g1:prop:principal-transport-laws}
	设 \(G\) 为有限维实 Lie 群，\(\pi:P\to M\) 为光滑主右 \(G\) 丛，\(M\) 无边界，\(\omega\) 为主联络。设 \(\gamma:[a,b]\to M\)、\(\eta:[c,d]\to M\) 均连续且有限分片光滑，并满足 \(\gamma(b)=\eta(c)\)。按\cref{g1:thm:principal-horizontal-lift}，以 \(\omega\) 水平提升的终点定义
	\[
		T_\gamma:P_{\gamma(a)}\longrightarrow P_{\gamma(b)}.
	\]
	它是 \(G\) 等变微分同胚，满足 \(T_\gamma(p g)=T_\gamma(p)g\)。以 \(\gamma*\eta\) 表示先走 \(\gamma\)、后走 \(\eta\) 的拼接，以 \(\bar\gamma(t)=\gamma(a+b-t)\) 表示反向路径。常值路径给出恒等映射，且有 \(T_{\gamma*\eta}=T_\eta\circ T_\gamma\)，而 \(T_{\bar\gamma}=T_\gamma^{-1}\)。若 \(h:[a',b']\to[a,b]\) 光滑非减，且 \(h(a')=a\)、\(h(b')=b\)，则 \(T_{\gamma\circ h}=T_\gamma\)，允许速度为零的停顿。
\end{proposition}
\begin{proof}
	若 \(q\) 水平，则 \(R_g\circ q\) 也水平，因为 \(R_g^*\omega=\operatorname{Ad}_{g^{-1}}\omega\)。其初值是 \(q(a)g\)，唯一性给出等变性。端点映射对初值光滑，由反向路径得到的逆映射也光滑。常值、拼接和反向的提升分别是常值提升、逐段连接的提升与反向提升；检验初值后使用唯一性即得前三条运算规律。最后，若时间换元为 \(t=h(\tau)\)，则 \(\omega((q\circ h)')=h'\omega(q')=0\)，没有除以 \(h'\)，所以停顿也不影响结论。
\end{proof}

对表示 \(\rho:G\to\operatorname{GL}(V)\)，关联丛中的向量 \([p,v]\) 沿 \(\gamma\) 变成 \([T_\gamma(p),v]\)。这正是\cref{g1:thm:associated-covariant-derivative}所构造联络的平行移动：在丛图中，\(q=s(\gamma)g\) 给出的向量分量为 \(\rho(g)v\)，对\eqref{g1:eq:principal-horizontal-ode}应用 \(\rho\) 后，其导数为 \(-\rho_*(A(\gamma'))\rho(g)v\)，恰好满足\eqref{g1:eq:parallel-transport-ode}。

\ProseParagraphBreak

若 \(\gamma\) 是以 \(x\) 为基点的环路，选定 \(p\in P_x\) 后唯一写成 \(T_\gamma(p)=p h_\gamma\)。运输与右作用等变，所以
\begin{equation}\label{g1:eq:principal-loop-composition}
	h_{\gamma*\eta}=h_\eta h_\gamma.
\end{equation}
在下一节将运输写成基本群表示时，这个次序正是需要取逆的原因。

\begin{example}[一般水平分布的提升可能中断]\label{g1:exm:incomplete-ehresmann-lift}
	在投影 \(\mathbb R^2\to\mathbb R\)、\((x,y)\mapsto x\) 上，令
	\[
		H_{(x,y)}=\operatorname{span}\{\partial_x+y^2\partial_y\}.
	\]
	它光滑且与竖直空间 \(\operatorname{span}\{\partial_y\}\) 互补，因而是一般纤维丛意义的水平分布。路径 \(x(t)=t\)、\(0\le t\le2\)，从 \((0,1)\) 出发的水平提升却须满足 \(y'=y^2\)，故 \(y(t)=(1-t)^{-1}\)，在 \(t=1\) 中断。纤维上的平移 \(y\mapsto y+c\) 不保持这个分布，所以它不是加法群主作用的主联络。由此可见，局部互补条件本身不能推出任意给定路径的完整提升；上述定理额外使用了主作用等变性。
\end{example}

\subsection{曲率结构方程}
\label{g1:subsec:1803:02}

若先沿一个水平方向、再沿另一个水平方向运动，与交换次序后的结果不同，差别可能含竖直方向。微分形式会把这种竖直偏差记录下来。对两个水平向量场 \(U,V\)，\eqref{g1:eq:exterior-one-form-evaluation}已经给出 \(\mathrm{d}\omega(U,V)=-\omega([U,V])\)。不过 \(\mathrm{d}\omega\) 在任意切向输入上未必只记录水平变化，还须补上一个 Lie 括号项。

\ProseParagraphBreak

为明确这个项的系数，选 \(\mathfrak g\) 的任一基 \(E_a\)，对 \(\mathfrak g\) 值形式 \(\alpha=\sum_a\alpha^aE_a\)、\(\beta=\sum_b\beta^bE_b\) 定义
\begin{equation}\label{g1:eq:lie-bracket-wedge}
	[\alpha\wedge\beta]\coloneqq\sum_{a,b}\alpha^a\wedge\beta^b[E_a,E_b].
\end{equation}
右边只是把双线性 Lie 括号作用于系数，因此与所选基无关。特别地，对一形式 \(\omega\)，有 \([\omega\wedge\omega](U,V)=2[\omega(U),\omega(V)]\)。这解释了下一定义中的 \(1/2\)。

\begin{definition}[主联络的曲率]\label{g1:def:principal-curvature}
	设 \(\omega\) 为光滑右主 \(G\) 丛 \(P\to M\) 上的联络形式。其\term[zhulianluoqulv]{曲率形式}[curvature form]定义为
	\begin{equation}\label{g1:eq:principal-curvature-structure}
		\Omega\coloneqq\mathrm{d}\omega+\tfrac12[\omega\wedge\omega]\in\Omega^2(P;\mathfrak g).
	\end{equation}
	若 \(\Omega=0\)，则称该主联络平坦。若 \(G\) 是矩阵 Lie 群，括号是矩阵交换子，上式也写成 \(\Omega=\mathrm{d}\omega+\omega\wedge\omega\)，此处后一个楔积同时使用矩阵乘法。
\end{definition}

\begin{proposition}[曲率的水平性与局部表达]\label{g1:prop:principal-curvature-properties}
	设 \(P\to M\) 为有限维实 Lie 群的光滑右主丛，\(\omega\) 为主联络，曲率为 \(\Omega\)。则
	\begin{enumerate}
		\item \(\Omega\) 是水平的，即任意一个输入为竖直向量时其值为零；并且 \(R_g^*\Omega=\operatorname{Ad}_{g^{-1}}\Omega\)。
		\item 对水平向量场 \(U,V\)，有 \(\Omega(U,V)=-\omega([U,V])\)。
		\item 若 \(s_i:U_i\to P\) 为局部截面，\(A_i=s_i^*\omega\)，则
		\begin{equation}\label{g1:eq:principal-local-curvature}
			F_i\coloneqq s_i^*\Omega=\mathrm{d}A_i+\tfrac12[A_i\wedge A_i],\qquad
			F_j=\operatorname{Ad}_{g_{ij}^{-1}}F_i\quad(s_j=s_i g_{ij}).
		\end{equation}
		\item 对光滑有限维线性表示 \(\rho:G\to\operatorname{GL}(V)\)，关联丛联络的局部曲率矩阵为 \(\rho_*(F_i)\)。
	\end{enumerate}
\end{proposition}
\begin{proof}
	对固定 \(X\in\mathfrak g\)，基本场 \(X^\#\) 的流为 \(R_{\exp(tX)}\)。对联络的等变公式求导，得到
	\[
		\mathcal L_{X^\#}\omega=-[X,\omega].
	\]
	\cref{g1:thm:cartan-formula}逐分量适用，而 \(\iota_{X^\#}\omega=X\) 为常值，故 \(\iota_{X^\#}\mathrm{d}\omega=-[X,\omega]\)。另一方面，\(\tfrac12\iota_{X^\#}[\omega\wedge\omega]=[X,\omega]\)。两项相消，得到 \(\iota_{X^\#}\Omega=0\)。每个竖直向量都是某个 \(X^\#\) 的点值，因此 \(\Omega\) 水平。固定群元素的伴随作用与外微分交换，又保持 Lie 括号；对结构方程取 \(R_g^*\) 便得到等变性。

	对水平 \(U,V\)，括号楔项为零，外微分的一形式公式化为 \(\Omega(U,V)=-\omega([U,V])\)，说明曲率由水平场括号的竖直分量确定。

	拉回与外微分交换，见\cref{g1:prop:exterior-pullback}，并且拉回保持\eqref{g1:eq:lie-bracket-wedge}，故 \(F_i\) 满足所列结构方程。计算 \(s_j=s_i g_{ij}\) 的微分时，一项是 \(\mathrm{d}R_{g_{ij}}\mathrm{d}s_i\)，另一项来自 \(g_{ij}\) 的变化，属于竖直空间。曲率的水平性消去后一项，等变性则将前一项变成 \(\operatorname{Ad}_{g_{ij}^{-1}}F_i\)。这正是曲率变换式中没有 \(\mathrm{d}g_{ij}\) 的原因。

	最后，\cref{g1:thm:associated-covariant-derivative}给出关联联络矩阵 \(\rho_*(A_i)\)。由\cref{g1:thm:curvature-structure-equation}，其曲率矩阵为
	\[
		\mathrm{d}(\rho_*A_i)+(\rho_*A_i)\wedge(\rho_*A_i)
		=\rho_*\bigl(\mathrm{d}A_i+\tfrac12[A_i\wedge A_i]\bigr),
	\]
	其中用了\cref{g1:prop:lie-homomorphism-differential}中 Lie 群同态的微分保持 Lie 括号这一结论。因此它等于 \(\rho_*(F_i)\)。
\end{proof}

曲率因而也可看作底空间上取值于伴随丛 \(\operatorname{ad}P=P\times_{\operatorname{Ad}}\mathfrak g\) 的二形式：在 \(x\) 处取 \(p\in P_x\)，把底空间的输入水平提升至 \(p\)，再取 \([p,\Omega_p(\cdot,\cdot)]\)。改取 \(pg\) 时，输入由 \(\mathrm{d}R_g\) 搬运，曲率值变为 \(\operatorname{Ad}_{g^{-1}}\Omega_p\)，恰好被关联丛的等价关系抵消。因此这个底空间对象无需选择局部截面；\(F_i\) 是它在所选截面中的表示。

\begin{example}[曲率与水平括号]\label{g1:exm:principal-additive-curvature}
	在 \(P=\mathbb R^2\times\mathbb R\to\mathbb R^2\) 上，让加法群 \(\mathbb R\) 作用于最后一个坐标 \(z\)，取 \(\omega=\mathrm{d}z+x\,\mathrm{d}y\)。这是主联络：它对 \(\partial_z\) 取值为 \(1\)，且在 \(z\) 平移下不变。水平分布由
	\[
		U=\partial_x,\qquad V=\partial_y-x\partial_z
	\]
	张成，而 \([U,V]=-\partial_z\)、\(\Omega=\mathrm{d}x\wedge\mathrm{d}y\)。于是 \(\Omega(U,V)=1=-\omega([U,V])\)。沿正向矩形边界提升时，\(z'=-x y'\)，走完后 \(z\) 减少矩形面积。这一竖直位移就是此处的环路返回，未必是角度或旋转。
\end{example}

\begin{theorem}[平坦联络的局部水平截面]\label{g1:thm:flat-principal-local-section}
	设 \(\pi:P\to M\) 为有限维实 Lie 群 \(G\) 的光滑主右丛，流形无边界，\(\omega\) 为主联络，\(\Omega\) 为其曲率；对任一局部截面 \(s\) 记 \(A_s=s^*\omega\)。以下条件等价：
	\begin{enumerate}
		\item \(\Omega=0\)。
		\item 水平分布 \(H=\ker\omega\) 对合。
		\item 每个 \(p\in P\) 都位于某个局部水平截面 \(s:U\to P\) 上，即 \(s(\pi(p))=p\) 且 \(s^*\omega=0\)。
	\end{enumerate}
	因此平坦主联络在足够小的截面中具有 \(A_s=0\) 的局部形式；这样的截面之间的过渡函数局部常值。
\end{theorem}
\begin{proof}
	曲率水平，故要判断它是否为零，只须代入两个水平输入。前一命题给出 \(\Omega(U,V)=-\omega([U,V])\)，所以它处处为零恰好等价于水平场的括号仍水平。这证明前两项等价。

	若 \(H\) 对合，\cref{g1:thm:frobenius}适用于总空间 \(P\) 上秩为 \(\dim M\) 的光滑分布，给出通过 \(p\) 的局部嵌入积分片 \(L\)。投影微分 \(\mathrm{d}\pi:H_p\to \mathrm{T}_{\pi(p)}M\) 是同构，所以\cref{g1:thm:inverse-function}使 \(\pi|_L\) 在 \(p\) 附近成为微分同胚。其局部逆 \(s\) 满足 \(\pi\circ s=\operatorname{id}\)，并且 \(\mathrm{d}s(\mathrm{T}U)\subset H\)，即为所求水平截面。反过来，若每个 \(p\) 都有这样的截面，那么 \(s^*\Omega=\mathrm{d}(s^*\omega)+\tfrac12[s^*\omega\wedge s^*\omega]=0\)。\(\mathrm{d}s\) 在该点张成全部水平空间，曲率的水平性遂说明 \(\Omega_p=0\)。

	若两个水平截面满足 \(s_j=s_i g_{ij}\)，联络变换式给出 \(0=\theta^L(g_{ij})\)，即 \((\mathrm{d}L_{g_{ij}^{-1}})\mathrm{d}g_{ij}=0\)。左平移微分可逆，所以 \(\mathrm{d}g_{ij}=0\)。光滑映射微分为零时，在一个充分小的连通坐标邻域中沿坐标线积分可知其值不变，故 \(g_{ij}\) 局部常值。
\end{proof}

这个结果是局部陈述。水平截面在绕行之后能否重新拼成同一个截面，还要看环路运输；下一节将保留这部分整体信息。

\subsection{Bianchi 恒等式的主丛形式}
\label{g1:subsec:1803:03}

曲率的局部表达在更换截面时按伴随表示变化。因此对曲率求导时，需要补上这一表示的联络项，才能将各局部结果粘成同一个对象。对水平且满足 \(R_g^*\eta=\operatorname{Ad}_{g^{-1}}\eta\) 的 \(\mathfrak g\) 值 \(k\) 形式，定义
\begin{equation}\label{g1:eq:principal-exterior-covariant-derivative}
	\mathrm{D}\eta\coloneqq\mathrm{d}\eta+[\omega\wedge\eta].
\end{equation}
这仍是水平等变形式。事实上，\(\mathcal L_{X^\#}\eta=-[X,\eta]\)，Cartan 公式与 \(\iota_{X^\#}\eta=0\) 给出 \(\iota_{X^\#}\mathrm{d}\eta=-[X,\eta]\)，而 \(\iota_{X^\#}[\omega\wedge\eta]=[X,\eta]\)，两项抵消。等变性由外微分及 Lie 括号与固定伴随变换相容得到。零次形式时插入项本来为零，同一计算也适用。局部截面中的公式为 \(s^*\mathrm{D}\eta=\mathrm{d}(s^*\eta)+[A\wedge s^*\eta]\)，正是伴随 Lie 代数丛联络的外协变微分。

\begin{theorem}[主联络的 Bianchi 恒等式]\label{g1:thm:principal-bianchi}
	设 \(\pi:P\to M\) 为有限维实 Lie 群 \(G\) 的光滑主右丛，\(\omega\) 为主联络，\(\Omega\) 为其曲率。对光滑局部截面 \(s:U\to P\)，记 \(A=s^*\omega\)、\(F=s^*\Omega\)。则
	\begin{equation}\label{g1:eq:principal-bianchi}
		\mathrm{D}\Omega=\mathrm{d}\Omega+[\omega\wedge\Omega]=0.
	\end{equation}
	在局部截面中，这成为 \(\mathrm{d}F+[A\wedge F]=0\)。矩阵群情形即 \(\mathrm{d}F+A\wedge F-F\wedge A=0\)。
\end{theorem}
\begin{proof}
	Lie 括号的 Jacobi 恒等式提供了所需的抵消。对次数为 \(p,q\) 的形式，由定义逐项交换标量楔积与 Lie 括号，得到
	\[
		[\alpha\wedge\beta]=-(-1)^{pq}[\beta\wedge\alpha],\qquad
		\mathrm{d}[\alpha\wedge\beta]=[\mathrm{d}\alpha\wedge\beta]+(-1)^p[\alpha\wedge\mathrm{d}\beta].
	\]
	特别地，\(\mathrm{d}[\omega\wedge\omega]=-2[\omega\wedge\mathrm{d}\omega]\)。另外，直接在三个切向输入 \(U,V,W\) 上展开有
	\begin{align*}
		[\omega\wedge[\omega\wedge\omega]](U,V,W)
		={}&2[\omega(U),[\omega(V),\omega(W)]]\\
		&+2[\omega(V),[\omega(W),\omega(U)]]\\
		&+2[\omega(W),[\omega(U),\omega(V)]]=0,
	\end{align*}
	最后一步就是 \(\mathfrak g\) 的 Jacobi 恒等式。因此结构方程给出
	\[
		\mathrm{d}\Omega=-[\omega\wedge\mathrm{d}\omega],\qquad
		[\omega\wedge\Omega]=[\omega\wedge\mathrm{d}\omega],
	\]
	相加即为零。拉回到局部截面得到局部公式。矩阵情形的一次形式与二次形式之间，括号楔积为 \([A\wedge F]=A\wedge F-F\wedge A\)，故与\eqref{g1:eq:bianchi-matrix}完全一致。
\end{proof}

这个恒等式适用于任意主丛联络。若 \(G\) 为交换群，括号项消失，恒等式退化为 \(\mathrm{d}\Omega=0\)；一般情形删去括号项就改变了所比较的对象。

\subsection{小环路和平行移动的变化}
\label{g1:subsec:1803:04}

在\cref{g1:prop:principal-curvature-properties}中，曲率记录了水平场括号的竖直部分。现在把它与实际运输联系起来：先计算一族路径变化时终点怎样变化，再说明一个小正向方环为何产生负的曲率面积项。两种计算都保持上一节的右主作用约定。

\begin{theorem}[水平运输的变分]\label{g1:thm:principal-transport-variation}
	设 \(G\) 为有限维实 Lie 群，\(\pi:P\to M\) 为带主联络 \(\omega\) 的光滑主右 \(G\) 丛，\(\Omega\) 为其曲率，\(\theta^L_g=(\mathrm{d}L_{g^{-1}})_g\) 为 \(G\) 的左 Maurer--Cartan 形式。设 \(\Gamma:I\times[a,b]\to M\) 为光滑路径族，\(I\subset\mathbb R\) 为开区间。设 \(c:I\to P\) 光滑，\(\pi(c(s))=\Gamma(s,a)\)。以 \(c(s)\) 为初值，沿 \(t\mapsto\Gamma(s,t)\) 的水平提升记为 \(q(s,t)\)。令
	\[
		W(s,t)=\omega_{q(s,t)}(\partial_s q).
	\]
	则
	\begin{equation}\label{g1:eq:principal-variation-intrinsic}
		W(s,b)=\omega_{c(s)}(c'(s))+\int_a^b\Omega_{q(s,t)}(\partial_tq,\partial_sq)\,\mathrm{d}t.
	\end{equation}
	若 \(\Gamma\) 的像在同一个截面 \(\sigma:U\to P\) 的定义域内，取 \(c(s)=\sigma(\Gamma(s,a))\)、\(q(s,t)=\sigma(\Gamma(s,t))u(s,t)\)，则 \(u(s,a)=e\)。记 \(A=\sigma^*\omega\)、\(F=\sigma^*\Omega\)、\(A_s(s,t)=A(\partial_s\Gamma)\)，有
	\begin{equation}\label{g1:eq:principal-variation-endpoints}
		\begin{split}
			\theta^L_{u(s,b)}(\partial_su(s,b))
			={}&A_s(s,a)-\operatorname{Ad}_{u(s,b)^{-1}}A_s(s,b)\\
			&+\int_a^b\operatorname{Ad}_{u(s,t)^{-1}}F_{\Gamma(s,t)}(\partial_t\Gamma,\partial_s\Gamma)\,\mathrm{d}t.
		\end{split}
	\end{equation}
\end{theorem}
\begin{proof}
	水平提升的参数依赖保证 \(q\) 光滑，因而可以对拉回形式求外微分。沿 \(t\) 方向水平意味着
	\[
		q^*\omega=W\,\mathrm{d}s.
	\]
	它与自身的括号楔积为零，所以曲率结构方程拉回后给出
	\[
		q^*\Omega=\mathrm{d}(W\,\mathrm{d}s)=\partial_tW\,\mathrm{d}t\wedge\mathrm{d}s.
	\]
	代入 \((\partial_t,\partial_s)\)，得到 \(\partial_tW=\Omega(q_t,q_s)\)。从 \(a\) 到 \(b\) 积分，并注意 \(q(s,a)=c(s)\)，便得第一式。积分中的值都在同一个有限维向量空间 \(\mathfrak g\) 内，按任一基逐分量积分即可。

	在所给局部截面中，联络平凡化公式给出
	\[
		W(s,t)=\operatorname{Ad}_{u(s,t)^{-1}}A_s(s,t)+\theta^L_{u(s,t)}(\partial_su(s,t)).
	\]
	曲率水平且等变，故 \(\Omega(q_t,q_s)=\operatorname{Ad}_{u^{-1}}F(\Gamma_t,\Gamma_s)\)。初端 \(u(s,a)=e\) 对 \(s\) 常值，因此 \(W(s,a)=A_s(s,a)\)。将这三项代入第一式，再把终端的 \(A_s\) 项移到右边，得到第二式。
\end{proof}

% !TeX root = ../Book1.tex
\begin{figure}[htbp]
	\centering
	\begin{tikzpicture}[font=\small]
		\fill[GeometryBlueLight] (0,0) rectangle (6,2.4);
		\draw[GeometryGray,dashed] (0,0)--(0,2.4);
		\draw[GeometryGray,dashed] (6,0)--(6,2.4);
		\draw[GeometryBlue,thick,->] (0,0)--(3.6,0);
		\draw[GeometryBlue,thick] (3.6,0)--(6,0);
		\draw[GeometryBlue,thick,->] (0,2.4)--(3.6,2.4);
		\draw[GeometryBlue,thick] (3.6,2.4)--(6,2.4);
		\draw[GeometryGray,->] (0,1.2)--(6,1.2);
		\draw[orange!80!black,thick,->] (2.6,0.4)--(2.6,1.95);
		\node[right,orange!80!black] at (2.6,1.75) {改变路径：\(\partial_s\)};
		\node[below] at (3,0) {沿路径前进：\(\partial_t\)};
		\node[above] at (3,2.4) {\(\Gamma(s_1,t)\)};
		\node[left,align=right] at (-0.1,1.2) {起点曲线\\\(\Gamma(s,a)\)};
		\node[right,align=left] at (6.1,1.2) {终点曲线\\\(\Gamma(s,b)\)};
		\node[below left] at (0,0) {\((s_0,a)\)};
		\node[below right] at (6,0) {\((s_0,b)\)};
	\end{tikzpicture}
	\caption{图画在路径族的参数域中，\(\Gamma\) 本身可以自交或降秩。比较相邻路径的运输时，中间窄带给出曲率项；两端的移动给出端点项。只有定端点变化才能直接省去这两项。}
	\label{g1:fig:principal-transport-variation}
\end{figure}


\cref{g1:fig:principal-transport-variation}把两种变化放在同一张参数图里：\(t\) 沿着路径前进，\(s\) 改变整条路径。公式中曲率的输入次序是 \((\Gamma_t,\Gamma_s)\)；改成 \((\Gamma_s,\Gamma_t)\) 就应补一个负号。若起点、终点随 \(s\) 移动，\eqref{g1:eq:principal-variation-endpoints}中的两个端点项不能省略。若两端固定，端点项消失，而积分中的 \(\operatorname{Ad}_{u^{-1}}\) 将不同位置的曲率放到提升所选的同一比较方式中；一般非交换群不能删去这一项。

\ProseParagraphBreak

特别地，对一族基点固定的环路，从同一个 \(p\in P_x\) 提升，写 \(q(s,b)=p h(s)\)，则
\[
	\theta^L_{h(s)}(h'(s))=\int_a^b\Omega(q_t,q_s)\,\mathrm{d}t.
\]
平坦时右侧为零，因而 \(h'(s)=0\)。这已经说明光滑定端点同伦不改变平坦运输；下一节将用局部水平截面把运输延伸到连续路径，并证明基本群所采用的连续同伦也保持它。

\begin{proposition}[小方环的曲率项]\label{g1:prop:principal-small-loop}
	设 \(\pi:P\to M\) 为光滑主右 \(G\) 丛，\(G\) 为有限维实 Lie 群，\(\mathfrak g=\operatorname{Lie}(G)\)，\(\omega\) 为给定主联络，\(\Omega\) 为其曲率。以 \(T_\gamma\) 表示 \(\omega\) 的水平运输。取内点 \(x_0\) 附近的坐标，使 \(x(x_0)=0\)，并取局部截面 \(\sigma\)，令 \(p=\sigma(x_0)\)、\(F=\sigma^*\Omega\)。假设 \(\dim M\ge2\)，并记 \(F_{12}(x_0)=F_{x_0}(\partial/\partial x^1,\partial/\partial x^2)\)。令 \(\gamma_\varepsilon\) 从原点沿第一坐标正向、第二坐标正向、第一坐标反向、第二坐标反向绕行边长 \(\varepsilon>0\) 的方环，其余坐标为零。写
	\[
		T_{\gamma_\varepsilon}(p)=p h_\varepsilon.
	\]
	当 \(\varepsilon\) 充分小时，\(h_\varepsilon\) 位于指数映射的局部可逆邻域，且
	\begin{equation}\label{g1:eq:principal-small-loop}
		\log h_\varepsilon=-\varepsilon^2F_{12}(x_0)+O(\varepsilon^3).
	\end{equation}
	余项用 \(\mathfrak g\) 上任一固定范数衡量。对矩阵 Lie 群，等价地有 \(h_\varepsilon=I-\varepsilon^2F_{12}(x_0)+O(\varepsilon^3)\)。
\end{proposition}
\begin{proof}
	先选一个在基点处一阶变化恰好抵消 \(A=\sigma^*\omega\) 的截面变换。令
	\[
		k(x)=\exp\left(-\sum_i x^iA_i(x_0)\right),\qquad \widetilde\sigma=\sigma k.
	\]
	由\cref{g1:prop:exponential-local-diffeomorphism}，\((\mathrm{d}\exp)_0=\operatorname{id}\)，且指数映射在零点附近有光滑逆。因为 \(k(x_0)=e\)，联络的局部变换公式给出 \(\widetilde A_i(x_0)=0\)。新旧截面在基点相同，故 \(p,h_\varepsilon\) 不变；曲率的等变变换又保证 \(\widetilde F(x_0)=F(x_0)\)。以下在这个新截面中计算，省去波浪号。

	把四条边都用坐标长度参数化，总时间为 \(4\varepsilon\)。沿全部边，系数 \(B(t)=A(\gamma_\varepsilon'(t))\) 分片光滑且 \(\|B(t)\|\le C\varepsilon\)，因为 \(A(x_0)=0\)。选单位元附近的指数坐标 \(\xi=\log g\)。在这张图中，右不变方程成为
	\[
		\xi'=-K(\xi)B(t),\qquad
		K(\xi)=(\mathrm{d}\log)_{\exp\xi}\circ(\mathrm{d}R_{\exp\xi})_e,
	\]
	其中 \(K\) 光滑且 \(K(0)=I\)。在固定小闭球中，\(K\) 有界且 \(K(\xi)=I+O(\|\xi\|)\)。从 \(\xi(0)=0\) 出发，在解尚未离开这个球时，积分估计给出 \(\|\xi(t)\|\le C'\varepsilon t\le4C'\varepsilon^2\)。若 \(\varepsilon\) 足够小，该界严格小于球半径，所以解不可能首次碰到球边界。于是全程 \(\xi=O(\varepsilon^2)\)，特别 \(h_\varepsilon\) 确实在 \(\log\) 的定义域内。

	将这个界代回微分方程，\(K(\xi)-I\) 对积分的贡献至多为 \(O(\varepsilon^2)O(\varepsilon)O(\varepsilon)=O(\varepsilon^4)\)，因此
	\[
		\log h_\varepsilon=-\oint_{\gamma_\varepsilon}A+O(\varepsilon^4).
	\]
	对 \(\mathfrak g\) 的各分量应用矩形 Stokes 公式\cref{g1:thm:stokes-corners}，再作一阶余项估计，得到
	\begin{align*}
		\oint_{\gamma_\varepsilon}A
		&=\int_0^\varepsilon\int_0^\varepsilon(\partial_1A_2-\partial_2A_1)(s,t,0,\ldots,0)\,\mathrm{d}s\,\mathrm{d}t\\
		&=\varepsilon^2F_{12}(x_0)+O(\varepsilon^3).
	\end{align*}
	最后一步用到了 \(A(x_0)=0\)，所以基点的交换子项消失。两式合并证明一般 Lie 群的结论。矩阵情形再用 \(\exp Z=I+Z+O(\|Z\|^2)\) 即得矩阵表达；前面的论证本身没有要求 \(G\) 具有忠实矩阵表示。
\end{proof}

这与\eqref{g1:eq:small-loop-curvature}的向量丛公式同号：负号来自水平方程 \(g'=-A(\gamma')g\) 及规定的正向绕行。曲率为零意味着所有这些面积首项消失，但单个环路恰好返回原点不意味着邻近曲率为零；而一般结构群非交换的情形也不能把面积首项直接改写成任意大小曲面的精确指数积分。

% 纲要条目排为编号小节标题，并进入目录与书签；说明段仍保留于纲要注释中。

% !TeX root = ../../Book1.tex
% 纲要标题：### 18.4 规范群与平坦模问题
% 纲要位置：几何学计划纲要.md，第 832 行。
% 正文状态：初稿；尚未正式审核。

% 纲要内容开始
% 
% * 规范变换及局部公式
% * 规范轨道和稳定子
% * 平坦丛的单值构造
% * 表示空间模共轭与平坦结构的对应范围
% * 圆周、环面和 Hopf 丛的对照
% 
% ---
% 
% 纲要内容结束

\section{规范群与平坦模问题}
\label{g1:sec:1804}

局部联络形式会随着截面变化，但联络还可以被一个真正作用在总空间上的自同构改变。怎样判断两份局部公式只是同一几何对象的不同表示，还是经由自同构成为等价的联络？平坦情形中，曲率已经全部消失，这个问题又留下什么可比较的数据？本节先分清这些操作，再用环路运输组织平坦结构的整体分类。

\subsection{规范变换及局部公式}
\label{g1:subsec:1804:01}

固定右主 \(G\) 丛 \(\pi:P\to M\)。这里 \(G\) 为有限维实 Lie 群，\(P,M\) 为有限维、无边界、Hausdorff、第二可数光滑流形。底点表示实际位置，纤维则容纳同一位置的不同标架或内部选择。保持底点、同时与右群作用相容的自同构，是本节要研究的变换。

\begin{definition}\label{g1:def:gauge-group}
	设 \(G\) 为有限维实 Lie 群，\(\pi:P\to M\) 为光滑主右 \(G\) 丛，\(P,M\) 无边界、Hausdorff、第二可数。若光滑微分同胚 \(\Phi:P\to P\) 满足 \(\pi\circ\Phi=\pi\) 和 \(\Phi(pg)=\Phi(p)g\)，则称其为 \(P\) 的一个\term[guifanbianhuan]{规范变换}[gauge transformation]。所有规范变换按函数复合组成\term[guifanqun]{规范群}[gauge group]，记为 \(\mathcal G(P)\)。\symidx{\(\mathcal G(P)\)：主丛的规范群}
\end{definition}

\begin{proposition}[规范变换的表示]\label{g1:prop:gauge-equivariant-functions}
	设 \(P\to M\) 为光滑右主 \(G\) 丛，\(G\) 为有限维实 Lie 群。规范变换与光滑映射 \(u:P\to G\) 一一对应，其对应式及条件为
	\begin{equation}\label{g1:eq:gauge-equivariant-function}
		\Phi_u(p)=p\,u(p),\qquad u(pg)=g^{-1}u(p)g.
	\end{equation}
	若局部截面满足 \(s_j=s_i g_{ij}\)，记 \(u_i=u\circ s_i\)，则
	\[
		u_j=g_{ij}^{-1}u_i g_{ij},\qquad
		\Phi_u(s_i(x)g)=s_i(x)u_i(x)g.
	\]
	在这些函数下，复合对应逐点乘法：\(\Phi_u\circ\Phi_v=\Phi_{uv}\)。
\end{proposition}
\begin{proof}
	每根纤维上的右作用自由且传递，所以 \(\Phi(p)\) 唯一写成 \(p u(p)\)。在局部乘积坐标 \(p=s_i(x)g\) 中，\(\Phi(s_i(x))=s_i(x)u_i(x)\) 给出光滑 \(u_i\)，等变性进一步给出 \(u(s_i(x)g)=g^{-1}u_i(x)g\)。故 \(u\) 光滑且满足所写条件。反之，满足该条件的光滑 \(u\) 给出覆盖恒等的等变光滑映射，其局部表达 \((x,g)\mapsto(x,u_i(x)g)\) 有光滑逆 \((x,g)\mapsto(x,u_i(x)^{-1}g)\)，因而为规范变换。

	令 \(p=s_i(x)\)、\(g=g_{ij}(x)\)，立即得到交叠公式。最后，
	\[
		\Phi_u(\Phi_v(p))=p\,v(p)\,u(pv(p))
		=p\,v(p)\bigl(v(p)^{-1}u(p)v(p)\bigr)=p\,u(p)v(p),
	\]
	证明复合规律。注意中间必须使用 \(u\) 的等变性，不能把 \(u(pv(p))\) 直接写成 \(u(p)\)。
\end{proof}

对共轭左作用 \(g\cdot a=gag^{-1}\)，\cref{g1:prop:associated-fiber-bundle}给出群纤维丛 \(\operatorname{Ad}(P)=P\times_GG\)。式\eqref{g1:eq:gauge-equivariant-function}说明规范群也可识别为这个丛的光滑截面群。大写 \(\operatorname{Ad}(P)\) 的纤维是群，小写 \(\operatorname{ad}(P)=P\times_G\mathfrak g\) 的纤维是 Lie 代数；联络之差取值于后者，不要把两个对象混为一谈。

\begin{proposition}[联络的规范变换]\label{g1:prop:gauge-connection-formula}
	设 \(G\) 为有限维实 Lie 群，\(\pi:P\to M\) 为光滑主右 \(G\) 丛，\(\omega\) 为主联络，\(\Omega\) 为其曲率。令 \(u:P\to G\) 光滑且满足 \(u(pg)=g^{-1}u(p)g\)，并定义规范变换 \(\Phi_u(p)=pu(p)\)。则 \(\Phi_u^*\omega\) 仍为主联络，并有
	\begin{equation}\label{g1:eq:gauge-global-connection}
		\Phi_u^*\omega=\operatorname{Ad}_{u^{-1}}\omega+u^*\theta^L,
	\end{equation}
	其中 \(\theta^L\) 是 \(G\) 的左 Maurer--Cartan 形式。取固定局部截面 \(s_i\)，令 \(u_i=u\circ s_i\)、\(A_i=s_i^*\omega\)、\(F_i=s_i^*\Omega\)，则变换后的局部形式和曲率满足
	\begin{equation}\label{g1:eq:gauge-local-connection}
		A_i^u=\operatorname{Ad}_{u_i^{-1}}A_i+u_i^*\theta^L,
		\qquad F_i^u=\operatorname{Ad}_{u_i^{-1}}F_i.
	\end{equation}
	矩阵群时，第一式为 \(A_i^u=u_i^{-1}A_i u_i+u_i^{-1}\mathrm{d}u_i\)。
\end{proposition}
\begin{proof}
	\(\Phi_u\) 与右作用交换，并把每个基本场 \(X^\#\) 推到同一个基本场，所以拉回保持联络形式的等变性和再现条件。由\cref{g1:thm:principal-connection-equivalence}，拉回确实是主联络。

	对过 \(p\) 的曲线 \(p(t)\)，微分曲线 \(p(t)u(p(t))\)。由乘积法则，其速度分成把 \(p'(0)\) 右移的部分，以及由 \(u(p(t))\) 的变化产生的竖直部分。联络在前者上取值 \(\operatorname{Ad}_{u(p)^{-1}}\omega_p(p'(0))\)，在后者上取值 \(\theta^L_{u(p)}((\mathrm{d}u)_p p'(0))\)。这正是\eqref{g1:eq:gauge-global-connection}；它也就是式\eqref{g1:eq:principal-connection-trivialization}中允许截面和群参数同时变化的计算。再沿 \(s_i\) 拉回，得到局部联络公式。

	曲率的定义与拉回交换，故新曲率为 \(\Phi_u^*\Omega\)。原曲率水平，刚才速度分解中的竖直部分对曲率求值没有贡献。保留下来的右移部分由等变性给出 \(\Phi_u^*\Omega=\operatorname{Ad}_{u^{-1}}\Omega\)，再沿截面拉回即得曲率公式。
\end{proof}

现在可以比较三种容易混淆的操作。被动地把截面 \(s_i\) 换成 \(s_i k_i\)，联络 \(\omega\) 不变，局部表达却按同一个公式变化；主动地施行 \(\Phi_u\)，则保持所选 \(s_i\) 不变，把联络改为 \(\Phi_u^*\omega\)。两者公式相同，是因为 \(\Phi_u\circ s_i=s_i u_i\)，但各片任取的换截面函数 \(k_i\) 不必满足 \(u_j=g_{ij}^{-1}u_i g_{ij}\)，所以并不自动拼成一个整体规范变换。

\ProseParagraphBreak

底空间坐标变化又是第三件事。若 \(x=x(y)\)，同一个一形式 \(A=\sum_iA_i\,\mathrm{d}x^i\) 在新坐标下写成 \(A'_a=\sum_iA_i\,\partial x^i/\partial y^a\)，只使用协变指标的 Jacobian 变换。若真正将底点按微分同胚 \(f:M\to M\) 移动，首先得到的是拉回丛 \(f^*P\) 及其拉回联络；还须给定 \(P\) 与 \(f^*P\) 的丛同构才能在同一个总空间上比较。这不属于本节覆盖 \(\operatorname{id}_M\) 的规范群。

\subsection{规范轨道和稳定子}
\label{g1:subsec:1804:02}

以 \(\mathcal A(P)\) 记 \(P\) 上全部光滑主联络。它由\cref{g1:prop:principal-connection-affine}成为以 \(\Omega^1(M;\operatorname{ad}(P))\) 为平移空间的仿射空间。令 \(\omega\cdot\Phi=\Phi^*\omega\)，这给出右规范作用，因为
\[
	(\omega\cdot\Phi)\cdot\Psi=\Psi^*\Phi^*\omega
	=(\Phi\circ\Psi)^*\omega=\omega\cdot(\Phi\circ\Psi).
\]
同一轨道中的联络称为规范等价。我们此时只定义等价类集合 \(\mathcal A(P)/\mathcal G(P)\)，尚未给这个无穷维商建立流形结构。

\ProseParagraphBreak

给定基点 \(x_0\) 和纤维点 \(p_0\in P_{x_0}\)，把沿分段光滑路径 \(\gamma\) 的水平运输记为 \(T_\gamma\)。\cref{g1:thm:principal-horizontal-lift,g1:prop:principal-transport-laws}保证它在整个紧参数区间存在，且与右作用交换。对基于 \(x_0\) 的环路，有唯一 \(h_\gamma\in G\) 满足 \(T_\gamma(p_0)=p_0h_\gamma\)。所有这样的 \(h_\gamma\) 构成 \(G\) 的子群，记为 \(\operatorname{Hol}_{p_0}(\omega)\)，称为该点的\term[zhulianluoholonomyqun]{环路平行移动群}[holonomy group of a principal connection]。它未必闭，这里不预设其具有某种嵌入 Lie 群结构。

\begin{theorem}[规范稳定子与环路运输]\label{g1:thm:gauge-stabilizer-centralizer}
	设 \(G\) 为有限维实 Lie 群，\(\pi:P\to M\) 为光滑主右 \(G\) 丛，\(M\) 非空、连通、有限维、无边界、Hausdorff、第二可数，\(\omega\) 为光滑主联络。固定 \(x_0\in M\) 及 \(p_0\in P_{x_0}\)。把规范变换写成 \(\Phi_u(p)=pu(p)\)，其中 \(u:P\to G\) 光滑且 \(u(pg)=g^{-1}u(p)g\)；令 \(\mathcal A(P)\) 为 \(P\) 上全部光滑主联络组成的空间。对基于 \(x_0\) 的连续有限分片光滑环路 \(\gamma\)，以 \(\omega\) 的水平运输 \(T_\gamma(p_0)=p_0h_\gamma\) 定义返回元素，令 \(\operatorname{Hol}_{p_0}(\omega)=\{h_\gamma\}\)。则取值映射 \(\Phi_u\mapsto u(p_0)\) 给出群同构
	\[
		\operatorname{Stab}_{\mathcal G(P)}(\omega)
		\cong Z_G(\operatorname{Hol}_{p_0}(\omega)),
	\]
	右侧为与全部环路运输元素交换的群元素所成的中心化子。特别地，固定 \(p_0\) 的规范子群在 \(\mathcal A(P)\) 上自由作用。
\end{theorem}
\begin{proof}
	先核对环路平行移动群确为子群。若先走 \(\gamma\) 再走 \(\eta\)，等变性给出
	\[
		T_{\gamma*\eta}(p_0)=T_\eta(p_0h_\gamma)
		=p_0h_\eta h_\gamma,
		\qquad h_{\gamma*\eta}=h_\eta h_\gamma.
	\]
	常路径给出 \(e\)，反向路径给出逆元；虽然所写乘法次序反转，全部取值仍对群乘法封闭。

	若 \(\Phi^*\omega=\omega\)，则 \(\Phi\) 保持水平分布，因此由提升唯一性与所有 \(T_\gamma\) 交换。记 \(\Phi(p_0)=p_0a\)。对环路，一方面 \(\Phi(T_\gamma p_0)=p_0a h_\gamma\)，另一方面 \(T_\gamma(\Phi p_0)=p_0h_\gamma a\)，故 \(a\) 与全部 \(h_\gamma\) 交换。对从 \(x_0\) 到 \(x\) 的任意分段光滑路径 \(\gamma\)，等变性又迫使
	\begin{equation}\label{g1:eq:gauge-stabilizer-transport}
		\Phi(T_\gamma(p_0)g)=T_\gamma(p_0)a g.
	\end{equation}
	连通光滑流形可用有限段坐标线连接任意两点，所以此式确定 \(\Phi\) 在整个 \(P\) 上的取值，证明取值映射单射。

	反过来，给定中心化子中的 \(a\)，用\eqref{g1:eq:gauge-stabilizer-transport}定义 \(\Phi\)。若另一路径 \(\beta\) 到达同一点，写 \(T_\beta(p_0)=T_\gamma(p_0)k\)。沿 \(\beta*\bar\gamma\) 运输可知 \(k\in\operatorname{Hol}_{p_0}(\omega)\)。同一点的两种纤维表示于是满足 \(g=kg'\)，而
	\[
		T_\beta(p_0)a g'=T_\gamma(p_0)ka g'
		=T_\gamma(p_0)a k g'=T_\gamma(p_0)a g.
	\]
	所以定义与路径无关。

	为核对光滑性，在任一点附近取坐标球，先沿一条固定路径到球心，再沿坐标径向线到各点。水平提升对参数光滑依赖，故 \(s(x)=T_{\gamma_x}(p_0)\) 是该球上的光滑截面。在此平凡化中 \(\Phi(s(x)g)=s(x)a g\)，所以 \(\Phi\) 光滑；把 \(a\) 换成 \(a^{-1}\) 给出光滑逆。定义也表明它与沿任意路径的水平运输交换，故保持水平分布，因而保持联络。最后，复合在 \(p_0\) 的取值是 \(ab\)，所以取值映射保持群运算。若 \(\Phi\) 还固定 \(p_0\)，则 \(a=e\)，上式给出 \(\Phi=\operatorname{id}\)，得到自由性结论。
\end{proof}

稳定子的出现是具体的对称现象。例如 \(G\) 的中心总包含在上述中心化子中；对交换群，只要底空间连通，保持联络的规范变换就由一个常群元素确定。使用带基点的规范群可以消去这些稳定子，但仅凭作用自由仍不能断言商是流形：还需要选择函数空间、拓扑和适当的局部分析。

\subsection{平坦丛的单值构造}
\label{g1:subsec:1804:03}

主联络 \(\omega\) 平坦，指 \(\Omega=0\)。\cref{g1:thm:flat-principal-local-section}已把这个条件转成局部存在水平截面 \(s\)，即 \(s^*\omega=0\)。如果两个这样的截面满足 \(s'=s k\)，局部变换公式给出 \(0=k^*\theta^L\)。\(\theta^L_g:\mathrm{T}_gG\to\mathfrak g\) 可逆，故 \(\mathrm{d}k=0\)，\(k\) 在每个交叠连通分支上为常数。这个局部常数性质使第十五章的平坦运输论证可以不依赖向量空间结构而继续成立。

\begin{proposition}[平坦主丛的同伦运输]\label{g1:prop:flat-principal-homotopy}
	设有限维实 Lie 群 \(G\) 的光滑右主丛 \(P\to M\) 带平坦主联络，\(M\) 为无边界、Hausdorff、第二可数光滑流形。水平运输唯一延拓到连续路径，且只依赖定端点连续同伦类。若 \(M\) 非空且单连通，则给定 \(p_0\in P_{x_0}\)，存在唯一整体水平截面 \(s\)，满足 \(s(x_0)=p_0\)；从而带联络的丛同构于具有乘积水平分布的 \(M\times G\)。
\end{proposition}
\begin{proof}
	在水平截面域 \(U\) 内定义
	\[
		T^U_{y\leftarrow x}(s(x)g)=s(y)g.
	\]
	若改用 \(s'=s k\)，而 \(x,y\) 在交叠的同一连通分支内，则 \(k(x)=k(y)\)，所以这个映射不变。连续路径的像可以细分成有限段，每段落在一个这样的截面域内；将局部运输依路径次序复合。增加分点不改变结果，因为 \(T^U_{z\leftarrow y}T^U_{y\leftarrow x}=T^U_{z\leftarrow x}\)。两种选择取共同细分后逐段比较，便得到定义无关性。对于分段光滑路径，水平截面坐标中的群分量恒定，故这里的运输与前节水平提升一致。

	对一个定端点连续同伦，将参数方形细分，使每个闭小格的像落在一个水平截面域内。这样的有限网格由紧性和\cref{g1:lem:lebesgue-number}得到。每个小格上，先横后竖和先竖后横都保持同一个群分量，因此边运输满足 \(v_{i+1,j}\circ h_{ij}=h_{i,j+1}\circ v_{ij}\)。将一行的这些等式按次序复合，得到 \(v_{mj}\circ H_j=H_{j+1}\circ v_{0j}\)。固定端点使左右两边的竖边运输都是恒等，故相邻行运输相同，最后得到同伦两端的运输相同。这就是\cref{g1:thm:flat-path-independence}证明中网格消去的同一个论证，未交换非交换群中的因子次序，也不需要先把连续同伦光滑化。

	若 \(M\) 单连通，从 \(x_0\) 到 \(x\) 的任意两条路径定端点同伦，故可定义 \(s(x)=T_\gamma(p_0)\)。在任一连通水平截面域 \(U\) 内，固定到其中一点的路径，再在 \(U\) 内追加路径，便知 \(s\) 等于一个局部水平截面右乘固定群元素。因此 \(s\) 光滑且水平。运输唯一性还保证通过 \(p_0\) 的水平截面唯一。映射 \((x,g)\mapsto s(x)g\) 是主丛微分同胚，并把保持 \(g\) 不变的乘积水平分布送到原水平分布。
\end{proof}

\begin{proposition}[主丛的单值表示]\label{g1:prop:flat-principal-monodromy}
	设 \(\pi:P\to M\) 为有限维实 Lie 群 \(G\) 的光滑主右丛，指定平坦光滑主联络 \(\omega\)，\(M\) 非空、连通、无边界、Hausdorff、第二可数。固定 \(x_0\in M\)、\(p_0\in P_{x_0}\)，对基于 \(x_0\) 的连续环路 \(\gamma\)，以按局部水平截面延伸到连续路径的平行移动 \(T_\gamma\) 定义 \(T_\gamma(p_0)=p_0h_\gamma\)。约定 \(\pi_1(M,x_0)\) 的乘积 \([\gamma][\eta]=[\gamma*\eta]\) 先走 \(\gamma\) 后走 \(\eta\)。则
	\begin{equation}\label{g1:eq:principal-monodromy-convention}
		\rho:\pi_1(M,x_0)\longrightarrow G,\qquad
		\rho([\gamma])=h_\gamma^{-1}
	\end{equation}
	是群同态。将 \(p_0\) 换成 \(p_0a\) 时，\(\rho\) 换成 \(a^{-1}\rho a\)。存在整体水平截面，当且仅当 \(\rho\) 平凡。
\end{proposition}
\begin{proof}
	前一命题使 \(h_\gamma\) 只依赖环路类。由于 \(h_{\gamma*\eta}=h_\eta h_\gamma\)，取逆得到 \(\rho([\gamma][\eta])=\rho([\gamma])\rho([\eta])\)。换基点后，\(T_\gamma(p_0a)=p_0h_\gamma a=(p_0a)(a^{-1}h_\gamma a)\)，故有共轭公式。

	整体水平截面使全部环路运输恒等，所以 \(\rho\) 平凡。反之，若所有环路运输恒等，则两条从 \(x_0\) 到 \(x\) 的路径 \(\gamma,\beta\) 满足 \(T_\beta^{-1}T_\gamma=T_{\gamma*\bar\beta}=\operatorname{id}\)。于是可沿路径定义 \(s(x)=T_\gamma(p_0)\)，并用前一命题中的局部公式证明它光滑且水平。
\end{proof}

这里取逆的原因与\cref{g1:prop:flat-monodromy}完全相同：道路乘积先走左边，运输复合却先作用右边。若 \(\tau:G\to\operatorname{GL}(V)\) 是表示，则关联向量丛的实际环路运输为 \(\tau(h_\gamma)\)，按第十五章约定得到的左单值表示恰是 \(\tau\circ\rho\)。

\begin{theorem}[由单值构造平坦主丛]\label{g1:thm:principal-monodromy-construction}
	设 \(M\) 为非空、连通、有限维、无边界、Hausdorff、第二可数光滑流形，\(G\) 为有限维实 Lie 群，\(\Gamma=\pi_1(M,x_0)\)。给定群同态 \(\rho:\Gamma\to G\)，令 \(\widetilde M\to M\) 为道路类万有覆盖，并取前置环路的左覆盖变换作用 \(L_\gamma\)。则
	\begin{equation}\label{g1:eq:flat-principal-representation-quotient}
		P_\rho=(\widetilde M\times G)/\Gamma,
		\qquad \gamma\cdot(q,g)=(L_\gamma q,\rho(\gamma)g)
	\end{equation}
	带有自然光滑右主 \(G\) 丛结构。局部保持第二因子恒定给出平坦主联络 \(\omega_\rho\)；相对于 \(p_0=[\widetilde x_0,e]\)，其单值恰为 \(\rho\)。
\end{theorem}
\begin{proof}
	流形的坐标球可缩，故\cref{g1:thm:path-class-universal-cover}适用。\cref{g1:prop:universal-deck-discontinuous}给出覆盖纤维上自由传递的左覆盖变换作用，满足 \(L_\gamma L_\delta=L_{\gamma\delta}\)。\(\rho\) 是同态，因此题中公式确为左作用。它与第二因子上的右 \(G\) 作用交换，所以 \([q,g]a=[q,ga]\) 良定义。

	取连通的均匀覆盖坐标球 \(U\)，选一张覆盖片的局部逆 \(\sigma_U:U\to\widetilde M\)。映射
	\[
		U\times G\longrightarrow P_\rho|_U,
		\qquad (x,g)\longmapsto[\sigma_U(x),g]
	\]
	是双射：覆盖纤维的任意点唯一写成 \(L_\gamma\sigma_U(x)\)，而 \([L_\gamma\sigma_U(x),g]=[\sigma_U(x),\rho(\gamma)^{-1}g]\)。它也是同胚，因为轨道商映射开，限制在一张覆盖片乘 \(G\) 后仍是开连续双射。

	若 \(\sigma_V=L_{c(x)}\sigma_U\)，则 \(c(x)\in\Gamma\) 局部常值：在一个更小的连通均匀覆盖邻域中，两局部逆各自落在固定覆盖片内，使它们在一点相等的覆盖变换便在整片上使它们相等。因而截面 \(s_U=[\sigma_U,e]\)、\(s_V=[\sigma_V,e]\) 满足 \(s_V=s_U\rho(c(x))^{-1}\)。过渡函数局部常值，尤其光滑。由这些乘积图得到光滑主丛；不同底点的分离由 \(M\) 的 Hausdorff 性给出，同一纤维的不同点在同一乘积图内分开。底空间有可数平凡化子覆盖，各乘积图均有可数基，故总空间第二可数。

	在每个局部截面中取 \(A_U=0\)。由于过渡函数局部常值，\cref{g1:thm:principal-connection-local-data}的粘合公式成立，故得到主联络，其曲率为零。沿环路 \(\gamma\) 的覆盖提升，从 \(\widetilde x_0\) 到 \(L_{[\gamma]}\widetilde x_0\)，保持 \(g=e\) 给出水平提升。因此
	\[
		T_\gamma(p_0)=[L_{[\gamma]}\widetilde x_0,e]
		=[\widetilde x_0,\rho([\gamma])^{-1}]
		=p_0\rho([\gamma])^{-1},
	\]
	按照\eqref{g1:eq:principal-monodromy-convention}取逆，恢复给定的 \(\rho\)。
\end{proof}

\subsection{表示空间模共轭与平坦结构的对应范围}
\label{g1:subsec:1804:04}

上述构造表明每个表示都能实现。为了得到分类，还必须证明任意平坦主丛都由自己的单值重建，并准确说明什么时候将两个对象视为相同。这里始终固定底流形及其恒等映射；允许底流形微分同胚的分类是另一问题。

\begin{theorem}[平坦主丛的分类]\label{g1:thm:flat-principal-classification}
	设 \(M\) 为非空、连通、有限维、无边界、Hausdorff、第二可数光滑流形，\(G\) 为有限维实 Lie 群，固定 \(x_0\in M\)。光滑右主 \(G\) 丛连同指定平坦主联络，在覆盖 \(\operatorname{id}_M\) 且保持联络的主丛同构下的等价类，与
	\[
		\operatorname{Hom}(\pi_1(M,x_0),G)/G
	\]
	一一对应，其中 \(G\) 按同时共轭作用。若还指定 \(p_0\in P_{x_0}\)，并要求同构保持该点，则对应未经共轭的表示集合 \(\operatorname{Hom}(\pi_1(M,x_0),G)\)。
\end{theorem}
\begin{proof}
	取带点平坦主丛 \((P,\omega,p_0)\)，令 \(\rho\) 为其单值。对 \(q=[\alpha]\in\widetilde M\)，定义
	\[
		\widehat F(q,g)=T_\alpha(p_0)g.
	\]
	同伦不变性保证它与 \(\alpha\) 的代表无关。若 \(\gamma\in\Gamma\)，则
	\[
		T_{\gamma*\alpha}(p_0)=T_\alpha(p_0)\rho(\gamma)^{-1},
		\qquad \widehat F(L_\gamma q,\rho(\gamma)g)=\widehat F(q,g).
	\]
	所以它下降为 \(F:P_\rho\to P\)。

	在一个足够小的连通水平截面域 \(U\) 内，先固定到其中一点的路径，再在 \(U\) 中追加路径，所得 \(T_\alpha(p_0)\) 给出一张水平截面。上述 \(F\) 在对应的局部乘积图中就是 \((x,g)\mapsto s(x)g\)。因此它在每根纤维上双射，有光滑局部逆，整体为主丛微分同胚；它还把乘积的水平曲线送到水平曲线，所以保持联络，并将 \([\widetilde x_0,e]\) 送到 \(p_0\)。这证明重建性。

	若两个带点平坦主丛同构，等变性及水平提升唯一性使它们的 \(h_\gamma\) 完全相同，故表示相同；反过来，相同表示使两者都等价于同一个 \(P_\rho\)，故带点等价类由表示分类。

	忘掉纤维点以后，可以把 \(p_0\) 换成任意 \(p_0a\)，表示正好变成 \(a^{-1}\rho a\)。每个这样的共轭都能由换点实现，且不改变底层带联络的主丛，故得到未指定纤维点的分类。也可直接核对：若 \(\rho'=a^{-1}\rho a\)，则 \([q,g]\mapsto[q,a^{-1}g]\) 给出 \(P_\rho\to P_{\rho'}\) 的保持联络的主丛同构。
\end{proof}

若一开始已固定底层光滑主丛 \(P\)，令 \(\mathcal A_{\mathrm{flat}}(P)\) 表示它的平坦联络集合，则正确的对应是
\begin{equation}\label{g1:eq:fixed-bundle-flat-moduli}
	\begin{split}
		&\mathcal A_{\mathrm{flat}}(P)/\mathcal G(P)\\
		&\quad\cong\{[\rho]\in\operatorname{Hom}(\pi_1(M,x_0),G)/G:
		P_\rho\cong P\text{ 作为光滑主丛}\}.
	\end{split}
\end{equation}
因为选取一个 \(P_\rho\to P\) 的光滑主丛同构可搬运联络，而换一个这样的同构恰好相差规范变换。右侧可能为空，也可能只占全部表示类的一部分。定理分类的是带联络对象，不是把全部表示类不加条件地塞进任意固定主丛。

\ProseParagraphBreak

“表示空间模共轭”在这里首先是一个集合商。即使给表示集合赋予逐点收敛拓扑，共轭商也未必为 Hausdorff 空间。例如取 \(M=S^1\)、\(G=\operatorname{GL}(2,\mathbb R)\)，表示由生成元的像决定。矩阵
\[
	U_t=\begin{pmatrix}1&t\\0&1\end{pmatrix}\qquad(t>0)
\]
都与 \(U_1\) 共轭，因为 \(U_t=\operatorname{diag}(t,1)U_1\operatorname{diag}(t,1)^{-1}\)，但 \(U_t\to I\)，而 \(I\) 不与 \(U_1\) 共轭。因此 \(U_1\) 的共轭轨道不是闭集，商中的对应单点也不闭。若希望得到光滑模空间，必须进一步讨论对象范围、稳定子与商拓扑，不能仅靠这份集合分类。

\subsection{圆周、环面和 Hopf 丛的对照}
\label{g1:subsec:1804:05}

先看圆周。\cref{g1:thm:circle-fundamental}给出 \(\pi_1(S^1)=\mathbb Z\)，所以一般 \(G\) 的平坦主丛由一个群元素的共轭类描述。一个元素描述的不是曲率，而是绕一周以后仍可能留下的整体运输。

\begin{example}[圆周上的相位]\label{g1:exm:circle-principal-flat-moduli}
	取 \(G=\operatorname{U}(1)\)，其 Lie 代数是 \(\mathrm{i}\mathbb R\)。在平凡主丛 \(S^1\times\operatorname{U}(1)\) 的标准截面中，取
	\[
		A_a=\mathrm{i}a\,\vartheta,\qquad a\in\mathbb R,
		\qquad \vartheta=(x\,\mathrm{d}y-y\,\mathrm{d}x)|_{S^1}.
	\]
	局部角度 \(\theta\) 中 \(\vartheta=\mathrm{d}\theta\)，但 \(\theta\) 不是圆周上的整体实值函数。一维底空间上曲率恒为零。沿正向一周，水平方程给出 \(h=\exp(-2\pi\mathrm{i}a)\)，所以 \(\rho(1)=\exp(2\pi\mathrm{i}a)\)。规范变换 \(u(e^{\mathrm{i}\theta})=e^{\mathrm{i}m\theta}\)、\(m\in\mathbb Z\)，使 \(A_a\) 变成 \(A_{a+m}\)。

	这些公式给出全部规范类。事实上，任意局部形式在整体标准截面中写成 \(A=\mathrm{i}f(\theta)\,\mathrm{d}\theta\)，其中 \(f\) 光滑且 \(2\pi\) 周期。令 \(a=(2\pi)^{-1}\int_0^{2\pi}f(t)\,\mathrm{d}t\)，则 \(\chi(\theta)=\int_0^\theta(a-f(t))\,\mathrm{d}t\) 周期，故 \(u=e^{\mathrm{i}\chi}\) 是整体规范变换，并将 \(A\) 变为 \(A_a\)。一般 \(u:S^1\to\operatorname{U}(1)\) 在实线提升为 \(e^{\mathrm{i}\chi(\theta)}\)，满足 \(\chi(\theta+2\pi)-\chi(\theta)=2\pi m\)；这是第五章圆周覆盖提升和整数绕数的结论，提升局部光滑。故 \(\int_{S^1}u^{-1}\mathrm{d}u=2\pi\mathrm{i}m\)，说明 \(A_a,A_b\) 规范等价当且仅当 \(b-a\in\mathbb Z\)。

	所以模集合为 \(\mathbb R/\mathbb Z\simeq\operatorname{U}(1)\)。每个表示都由上述平凡丛上的联络实现，分类定理于是还说明，圆周上每个平坦主 \(\operatorname{U}(1)\) 丛的底层光滑丛都平凡。然而只有 \(a\in\mathbb Z\) 的类具有整体水平截面。
\end{example}

\begin{example}[环面上的两个相位]\label{g1:exm:torus-principal-flat-moduli}
	由\cref{g1:prop:product-fundamental}，\(\pi_1(T^2)\cong\mathbb Z^2\)。一般 \(G\) 的表示由一对交换元素 \((a,b)\) 决定，等价关系为同时共轭。交换条件来自环面的基本群关系，不能省略成任意两个群元素。

	对 \(\operatorname{U}(1)\)，在 \(T^2\times\operatorname{U}(1)\) 上取
	\[
		A=\mathrm{i}(a\,\mathrm{d}\theta+b\,\mathrm{d}\varphi),\qquad a,b\in\mathbb R,
	\]
	其中两角度微分分别是两个圆周因子的整体角度一形式。曲率为零，两个生成元的单值为 \((e^{2\pi\mathrm{i}a},e^{2\pi\mathrm{i}b})\)。每一对相位都能这样得到，所以分类定理表明全部平坦主 \(\operatorname{U}(1)\) 丛均在这份平凡底层丛上实现；其规范类为 \((\mathbb R/\mathbb Z)^2\)。具体的 \(u=e^{\mathrm{i}(m\theta+n\varphi)}\) 给出整数平移 \((a,b)\mapsto(a+m,b+n)\)，而不同相位不能规范等价，因为交换群中的共轭没有作用。规范稳定子在每一类上均为常数 \(\operatorname{U}(1)\)。
\end{example}

前两例中，局部平坦性与非平凡整体运输可以并存。下面的圆周主丛则连平坦联络也不存在；我们用一个直接积分证明这一点，不把障碍留给后卷的示性类理论。

\begin{example}[Hopf 丛与曲率通量]\label{g1:exm:hopf-principal-curvature}
	令 \(\operatorname{U}(1)\) 在 \(S^3\subset\mathbb C^2\) 上右作用为 \((z_0,z_1)\lambda=(z_0\lambda,z_1\lambda)\)。作用自由且固有，商为 \(\mathbb{CP}^1\)，得到\term[hopfcong]{Hopf 丛}[Hopf bundle]
	\[
		\operatorname{U}(1)\longrightarrow S^3\longrightarrow\mathbb{CP}^1.
	\]
	这里 \(\mathbb{CP}^1\) 微分同胚于 \(S^2\)：在图 \([1:w]\)、\(w=x+\mathrm{i}y\) 中，映射为
	\[
		[1:w]\longmapsto\frac{(2x,2y,1-|w|^2)}{1+|w|^2}.
	\]
	其像是去掉南极的单位球面，逆由 \(w=(X+\mathrm{i}Y)/(1+Z)\) 给出；在另一图 \([v:1]\) 中代入 \(w=1/v\)，便在南极光滑延拓，且给出相应光滑局部逆。这也说明可用两张射影图描述球面。下面取复坐标 \((x,y)\) 给出的定向。

	总空间上的 \(\mathrm{i}\mathbb R\) 值一形式
	\[
		\omega=\overline z_0\,\mathrm{d}z_0+\overline z_1\,\mathrm{d}z_1
	\]
	是主联络。它在 \(S^3\) 上纯虚，因为实部为 \(\tfrac12\mathrm{d}(|z_0|^2+|z_1|^2)=0\)；它在基本向量 \((z_0\mathrm{i}t,z_1\mathrm{i}t)\) 上取值 \(\mathrm{i}t\)，且在右乘固定 \(\lambda\in\operatorname{U}(1)\) 下不变。再现条件与等变性因而都成立。

	在 \(U_0=\{z_0\ne0\}\)、\(U_1=\{z_1\ne0\}\) 上分别取截面
	\[
		s_0(w)=\frac{(1,w)}{\sqrt{1+|w|^2}},\qquad
		s_1(v)=\frac{(v,1)}{\sqrt{1+|v|^2}}.
	\]
	交叠处 \(v=1/w\)，直接相乘得到 \(s_1(1/w)=s_0(w)g_{01}(w)\)，其中
	\begin{equation}\label{g1:eq:hopf-transition}
		g_{01}(w)=\frac{|w|}{w}.
	\end{equation}
	因此局部截面不是任意两份独立坐标：在重叠环带上，它们相差一个沿圆周绕行的群元素。

	记 \(r=1+|w|^2\)。将 \(s_0\) 代入 \(\omega\)，微分归一化因子后得到
	\[
		A_0=\frac{\overline w\,\mathrm{d}w-w\,\mathrm{d}\overline w}{2r}
		=\frac{\mathrm{i}(x\,\mathrm{d}y-y\,\mathrm{d}x)}{1+x^2+y^2}.
	\]
	因为群为交换群，曲率在底空间上为整体二形式 \(F\)，其此图表达为
	\begin{equation}\label{g1:eq:hopf-curvature-flux}
		\begin{gathered}
			F|_{U_0}=\mathrm{d}A_0
			=\frac{\mathrm{d}\overline w\wedge\mathrm{d}w}{(1+|w|^2)^2}
			=\frac{2\mathrm{i}\,\mathrm{d}x\wedge\mathrm{d}y}{(1+x^2+y^2)^2},\\
			\frac1{2\pi\mathrm{i}}\int_{\mathbb{CP}^1}F=1.
		\end{gathered}
	\end{equation}
	最后一个积分可在此图中用极坐标计算：去掉的单点对光滑二形式积分没有贡献，而 \(\int_0^\infty 2R(1+R^2)^{-2}\,\mathrm{d}R=1\)，角变量积分为 \(2\pi\)。曲率在另一图光滑，亦保证这一非紧图积分就是整个紧曲面的积分。

	若该主丛有整体截面 \(s\)，则整体一形式 \(A=s^*\omega\) 满足 \(F=\mathrm{d}A\)。由闭曲面上的 Stokes 定理，\(\int F=0\)，与\eqref{g1:eq:hopf-curvature-flux}矛盾，所以丛不平凡。更强地，任何另一主联络与 \(\omega\) 的差都为 \(\pi^*\beta\)，其中 \(\beta\) 是底空间上的 \(\mathrm{i}\mathbb R\) 值一形式；这是\cref{g1:prop:principal-connection-affine}在交换群、平凡伴随表示下的结论。其曲率为 \(F+\mathrm{d}\beta\)，积分仍为 \(2\pi\mathrm{i}\)，所以不可能恒为零。Hopf 丛没有平坦主联络。
\end{example}

\begin{center}
	\small
	\begin{tabular}{p{.19\textwidth}p{.32\textwidth}p{.37\textwidth}}
		\toprule
		底层主丛 & 平坦联络与规范类 & 整体信息 \\
		\midrule
		\(S^1\times\operatorname{U}(1)\) & 相位 \(e^{2\pi\mathrm{i}a}\) & 曲率为零仍可有非平凡环路运输。 \\
		\(T^2\times\operatorname{U}(1)\) & 两个相位 & 两个基本圆周各自记录单值。 \\
		Hopf 丛 \(S^3\to S^2\) & 不存在平坦联络 & 曲率积分非零，且不随联络改变。 \\
		\bottomrule
	\end{tabular}
\end{center}

表中的区别同时涉及联络和底层丛。圆周与环面的例子表明，平凡丛不要求所选平坦联络有整体水平截面；Hopf 例子则说明，有些底层丛连平坦联络都不容许。另一方面，\(S^2\) 单连通，\cref{g1:prop:sphere-simply-connected}与分类定理也说明：其上所有平坦主丛都必须平凡。这与刚才的直接曲率积分相互核对，但不能用“局部可以令联络为零”来跳过整体障碍。


\begin{chaptersummary}
	主丛的局部截面满足 \(s_j=s_i g_{ij}\)，三重交叠给出 \(g_{ij}g_{jk}=g_{ik}\)；反过来，光滑余循环通过局部乘积的粘合重建主丛。表示把主丛变成关联向量丛，其截面对应总空间上的等变函数。对闭子群 \(H\subseteq G\)，结构群约化等价于在 \(P/H\to M\) 中选取光滑截面。

	\ProseParagraphBreak

	主联络既是右平移不变的垂直补空间，也是满足再现与伴随等变条件的一形式 \(\omega\)。其局部形式 \(A_i=s_i^*\omega\) 按 \(A_j=\operatorname{Ad}_{g_{ij}^{-1}}A_i+g_{ij}^*\theta^L\) 变换。单位分解保证整体存在，两联络之差属于 \(\Omega^1(M;\operatorname{ad}(P))\)。表示 \(\rho\) 诱导局部协变微分 \(\mathrm{d}+\rho_*(A_i)\)；标架丛给出与线性联络的双向对应，一般关联表示则可能丢失主联络信息。

	\ProseParagraphBreak

	主作用的等变性使水平提升在给定路径的整个紧参数区间存在，无需群紧。曲率 \(\Omega=\mathrm{d}\omega+\tfrac12[\omega\wedge\omega]\) 是水平等变二形式，记录水平场括号的垂直部分，并满足 \(\mathrm{D}\Omega=0\)。它的积分控制运输随路径的变化；正向小方环满足 \(\log h_\varepsilon=-\varepsilon^2F_{12}+O(\varepsilon^3)\)。平坦性等价于局部存在水平截面，其交叠过渡函数局部常值。

	\ProseParagraphBreak

	规范变换是覆盖底空间恒等映射的主丛自同构。它可以改变联络，而换局部截面只改变表达，底空间坐标变化则作用于微分形式的坐标指标。连通底空间上，保持联络的规范变换由一个纤维点处的取值确定，这些取值组成全部环路运输的中心化子。规范商因而可能有稳定子，不能仅从等价关系就认定它具有流形结构。

	\ProseParagraphBreak

	平坦运输只依赖定端点同伦类。若 \(T_\gamma(p_0)=p_0h_\gamma\)，则 \(\rho([\gamma])=h_\gamma^{-1}\) 按本书道路乘法约定成为同态。在非空连通底空间上，将平坦主丛连同其联络按保持联络且覆盖恒等映射的同构视为等价，则其等价类由 \(\operatorname{Hom}(\pi_1(M),G)/G\) 分类；预先固定底层主丛时，只保留能够重建该丛的表示类。圆周和环面的平坦圆周主丛分别记录一个与两个相位，曲率为零仍可没有整体水平截面。Hopf 丛的曲率积分在本章定向下为 \(2\pi\mathrm{i}\)，不随主联络改变，因而没有平坦联络。
\end{chaptersummary}

\BookExercises

以下习题依次对应主丛构造、联络的等价描述、水平运输与曲率、规范等价与平坦结构。计算时固定右作用及局部截面的变换方向；涉及环路时，分别写出实际运输元素与按本书约定取逆得到的单值表示。

\begin{exercise}\label{g1:ex:principal-pullback}
	设 \(\pi:P\to M\) 为有限维实 Lie 群 \(G\) 的光滑主丛，\(f:N\to M\) 为光滑映射，所有流形均有限维、无边界、Hausdorff、第二可数。
	\begin{enumerate}
		\item 证明 \(f^*P=\{(x,p)\in N\times P:f(x)=\pi(p)\}\) 在子空间拓扑下自然成为主 \(G\) 丛，投影为 \((x,p)\mapsto x\)，作用为 \((x,p)g=(x,pg)\)。求其过渡函数。
		\item 若 \(Q\to N\) 为主 \(G\) 丛，\(F:Q\to P\) 是覆盖 \(f\) 的光滑 \(G\) 等变映射，证明存在唯一覆盖 \(\operatorname{id}_N\) 的主丛同构 \(\widehat F:Q\to f^*P\)，使第二因子投影与 \(\widehat F\) 的复合为 \(F\)。
	\end{enumerate}
\end{exercise}
\begin{solution}
	1. 设 \(P|_{U_i}\) 的局部截面为 \(s_i\)。在 \(f^{-1}(U_i)\) 上定义
	\[
		(x,g)\longmapsto(x,s_i(f(x))g).
	\]
	它到相应部分的 \(f^*P\) 双射。在 \(N\times P|_{U_i}\) 的乘积坐标 \((x,y,g)\) 中，\(f^*P\) 是由 \(y=f(x)\) 描述的图；\((x,g)\) 为光滑子流形坐标，并与子空间拓扑相同。因此 \(f^*P\) 为嵌入子流形，继承 Hausdorff 性和第二可数性。上述图中作用是 \((x,g)a=(x,ga)\)，故它是主丛。若 \(s_j=s_i g_{ij}\)，拉回截面 \(\widetilde s_i(x)=(x,s_i(f(x)))\) 满足
	\[
		\widetilde s_j(x)=\widetilde s_i(x)g_{ij}(f(x)),
	\]
	所以过渡函数是 \(g_{ij}\circ f\)，定义在 \(f^{-1}(U_i\cap U_j)\) 上。

	2. 设 \(\pi_Q:Q\to N\)，唯一可能的公式是
	\[
		\widehat F(q)=(\pi_Q(q),F(q)).
	\]
	覆盖条件 \(\pi(F(q))=f(\pi_Q(q))\) 保证其值属于拉回丛。映射光滑、等变，并覆盖恒等映射。对每个 \(x\in N\)，两边纤维都是自由传递的 \(G\) 集；固定 \(q_0\in Q_x\) 后，\(\widehat F(q_0g)=\widehat F(q_0)g\)，所以逐纤维双射。在 \(Q\) 的局部截面 \(t\) 下，\(F(t(x))\) 给出拉回丛的局部截面，\(\widehat F\) 在这两组平凡化中是 \((x,g)\mapsto(x,g)\)。故逆也光滑，从而为主丛同构。第一分量由覆盖恒等映射确定，第二分量由复合条件确定，所以唯一。
\end{solution}

\begin{exercise}\label{g1:ex:mobius-principal-double-cover}
	在 \(P=S^1\subset\mathbb C\) 上令 \(\operatorname{O}(1)=\{1,-1\}\) 右作用为乘法，投影 \(\pi:P\to S^1\) 为 \(\pi(z)=z^2\)。
	\begin{enumerate}
		\item 证明这是主 \(\operatorname{O}(1)\) 丛，且没有连续全局截面。
		\item 对符号表示 \(\rho(\varepsilon)v=\varepsilon v\)（\(v\in\mathbb R\)），证明 \(P\times_\rho\mathbb R\) 同构于 Möbius 线丛 \(L=(\mathbb R\times\mathbb R)/((t+2\pi,v)\sim(t,-v))\)。
		\item 对平凡一维表示 \(\rho_0(\varepsilon)=1\)，证明关联线丛平凡。说明这一比较为何不能仅从一个关联丛恢复主丛。
	\end{enumerate}
\end{exercise}
\begin{solution}
	1. 每个短圆弧上可选择光滑角度 \(t(x)\)，则 \(s(x)=e^{\mathrm{i}t(x)/2}\) 是 \(\pi\) 的局部截面。每个平方根唯一写成 \(s(x)\varepsilon\)，其中 \(\varepsilon=\pm1\)，所以得到等变平凡化。若存在连续全局截面 \(s\)，映射 \((x,\varepsilon)\mapsto s(x)\varepsilon\) 就是从 \(S^1\times\{1,-1\}\) 到 \(S^1\) 的连续双射。源紧、目标 Hausdorff，故它是同胚；但源有两个非空的开闭分支，目标连通，矛盾。因此连连续截面也不存在。

	2. 定义
	\[
		[e^{\mathrm{i}t/2},v]\longmapsto[t,v]\in L.
	\]
	同一个 \(z=e^{\mathrm{i}t/2}\) 的参数可改变 \(4\pi k\)，而 \([t+4\pi k,v]=[t,v]\)，故与角度选择无关。关联关系把 \((z,v)\) 与 \((-z,-v)\) 视为相同，后者可用 \((t+2\pi,-v)\) 表示；在 \(L\) 中 \([t+2\pi,-v]=[t,v]\)，所以公式与关联商关系相容。反向公式为 \([t,v]\mapsto[e^{\mathrm{i}t/2},v]\)，同样良定义。两式互逆，且在任意短圆弧的角度坐标中都为恒等纤维坐标，故是光滑线丛同构。\cref{g1:exm:mobius-line-bundle}证明的反周期条件 \(a(t+2\pi)=-a(t)\) 还说明，该关联丛没有处处非零截面。

	3. 平凡表示下，\((z,v)\sim(-z,v)\)，所以 \([z,v]\mapsto(z^2,v)\) 是到 \(S^1\times\mathbb R\) 的线丛同构；局部平方根给出光滑逆。主丛 \(P\) 仍是第 1 问中非平凡的双覆盖，而任意主 \(\operatorname{O}(1)\) 丛用平凡表示关联得到的线丛都具有同一形式 \(M\times\mathbb R\)。因此表示可以把主丛的区别全部消去；关联构造的输出必须连同所用表示一起理解。
\end{solution}

\begin{exercise}\label{g1:ex:associated-tensor-operations}
	设 \(P\to M\) 为有限维实 Lie 群 \(G\) 的光滑主丛，\(\rho:G\to\operatorname{GL}(V)\)、\(\tau:G\to\operatorname{GL}(W)\) 为同一基域 \(\mathbb K=\mathbb R\) 或 \(\mathbb C\) 上的有限维光滑表示。记 \(E=P\times_\rho V\)、\(F=P\times_\tau W\)。证明存在自然向量丛同构
	\[
		P\times_{\rho\otimes\tau}(V\otimes W)\cong E\otimes F,
		\qquad P\times_{\rho^*}V^*\cong E^*,
	\]
	其中 \(\rho^*(g)\lambda=\lambda\circ\rho(g^{-1})\)。写出同构公式，并说明两条不同关联丛的向量为什么要用同一个 \(p\in P_x\) 表示后，才能直接对分量作张量积。
\end{exercise}
\begin{solution}
	固定 \(p\in P_x\)，分别有线性同构 \(V\to E_x\)、\(W\to F_x\)。它们的张量积给出
	\[
		[p,v\otimes w]\longmapsto[p,v]\otimes[p,w],
	\]
	再按线性延拓到任意张量。将 \(p\) 换为 \(pg\) 时，同一个左侧元素的分量变成 \(\rho(g^{-1})v\otimes\tau(g^{-1})w\)，右侧两个类分别不变，故公式良定义。固定 \(p\) 时它是两个线性同构的张量积，因此逐纤维可逆；局部取 \(p=s_i(x)\)，公式是固定向量空间中的张量积识别，故它及逆光滑。

	对偶情形，把 \([p,\lambda]\) 送到 \(E_x\) 上的线性泛函
	\[
		[p,v]\longmapsto\lambda(v).
	\]
	若使用 \(pg\)，同一个对偶类由 \(\rho^*(g^{-1})\lambda=\lambda\circ\rho(g)\) 表示，同一个向量由 \(\rho(g^{-1})v\) 表示，所以求值仍为 \(\lambda(\rho(g)\rho(g^{-1})v)=\lambda(v)\)。因此公式独立于参照；它逐纤维为对偶空间的线性同构，在局部截面中是固定的求值配对，故光滑可逆。

	若两个向量最初表示为 \([p,v]\)、\([q,w]\)，其中 \(q=pg\)，则第二个向量应先改写为 \([p,\tau(g)w]\)。因此它们的张量积对应 \([p,v\otimes\tau(g)w]\)，不能直接使用 \(v\otimes w\)。两列数字所属的参照不同；共同参照正是把两个纤维对象的张量积正确翻译为分量运算所需的数据。
\end{solution}

\begin{exercise}\label{g1:ex:orientation-reduction-cover}
	设 \(E\to M\) 为秩 \(r\geq1\) 的实光滑向量丛，\(M\) 有限维、无边界、Hausdorff、第二可数。令 \(P=\operatorname{Fr}(E)\)、\(H=\operatorname{GL}^+(r,\mathbb R)\)。证明 \(P/H\to M\) 是纤维有两个点的光滑覆盖，并将其纤维解释为 \(E_x\) 的两种定向。证明 \(E\) 可定向当且仅当该覆盖有光滑截面，等价地，当且仅当该覆盖同构于 \(M\times\{+1,-1\}\)。若 \(M\) 非空且连通，证明可定向时恰有两种定向。
\end{exercise}
\begin{solution}
	行列式符号给出 \(\operatorname{GL}(r,\mathbb R)/H\cong\{+1,-1\}\)。\(H\) 既开又闭，因而是闭 Lie 子群；\cref{g1:thm:reduction-section-criterion}给出 \(P/H\) 的局部乘积 \(U_i\times\{+1,-1\}\)。每个 \(U_i\) 的逆像分成两个开集，各自微分同胚地投影到 \(U_i\)，所以这是光滑双覆盖。同一纤维中两个标架相差正行列式矩阵，当且仅当它们给出相同定向，因此两个陪集就是两种定向。过渡作用由 \(\operatorname{sgn}\det g_{ij}\) 决定。

	光滑截面在每点连续地选择一个定向，正是向量丛可定向的含义。在局部双覆盖中，连续截面的第二分量取值于离散两点集，因而局部常值，所以它自动光滑。这也表明连续定向与光滑定向在这里没有额外差别。若给定一个截面 \(\sigma\)，把 \((x,+1)\) 送到 \(\sigma(x)\)，把 \((x,-1)\) 送到同一纤维中另一点，局部是两个片的排列，故给出覆盖同构。反过来，乘积覆盖的常值正号截面经同构给出截面。

	当 \(M\) 连通时，任何另一个定向与指定定向的比较符号是连续函数 \(M\to\{+1,-1\}\)，只能为常值。两个常值都能实现，所以恰有两种定向。非空条件排除空底空间只有一个空选择的情形；连通条件则排除不同连通分支可以独立反转定向的情形。
\end{solution}

\begin{exercise}\label{g1:ex:principal-determinant-connection}
	设 \(E\to M\) 为秩 \(r\geq1\) 的实光滑向量丛，\(\nabla\) 在局部标架 \(e=(e_1,\ldots,e_r)\) 中的联络矩阵为 \(A\)。把行列式线丛 \(\det E=\Lambda^rE\) 看成标架丛经表示 \(\det:\operatorname{GL}(r,\mathbb R)\to\mathbb R^\times\) 得到的关联线丛。
	\begin{enumerate}
		\item 求其诱导联络在 \(e_1\wedge\cdots\wedge e_r\) 中的矩阵，并核对与逐因子求协变导数的结果一致。
		\item 两个 \(E\) 上的联络在 \(\det E\) 上诱导相同联络，当且仅当它们的差满足什么条件？在秩 \(2\) 的平凡丛上给出不同联络诱导相同联络的例子。
	\end{enumerate}
\end{exercise}
\begin{solution}
	1. 在单位矩阵处，行列式的微分为 \(X\mapsto\operatorname{tr}X\)。这可由多重线性直接计算：\(\det(I+tX)\) 的一次项必须从恰好一列取 \(tX\)，其余列取标准基，非零贡献恰为 \(t\sum_aX^a{}_a\)。因此由\cref{g1:thm:associated-covariant-derivative}，行列式线的局部联络形式为 \(\operatorname{tr}A\)。逐因子求导也给出
	\[
		\nabla(e_1\wedge\cdots\wedge e_r)
		=\sum_{b=1}^r e_1\wedge\cdots\wedge(\nabla e_b)\wedge\cdots\wedge e_r
		=(\operatorname{tr}A)(e_1\wedge\cdots\wedge e_r),
	\]
	因为 \(\nabla e_b=\sum_a A^a{}_b e_a\)，而 \(a\ne b\) 的项有重复因子，楔积为零。

	2. 记两联络之差为 \(B\in\Omega^1(M;\operatorname{End}E)\)。所得行列式线联络之差为 \(\operatorname{tr}B\)，所以条件是 \(\operatorname{tr}B(v)=0\) 对所有底空间切向量 \(v\) 成立。这是整体条件，因为换标架时 \(B\) 共轭变换，而迹不变。在 \(\mathbb R\times\mathbb R^2\) 的标准标架中取
	\[
		A_0=0,\qquad A_1=\begin{pmatrix}1&0\\0&-1\end{pmatrix}\mathrm{d}x.
	\]
	两者对常截面 \(e_1\) 的协变导数分别为零和 \(e_1\mathrm{d}x\)，故不同；但迹均为零，所以在行列式线上都诱导标准联络。行列式记录总体积变化，不能记录这类互相抵消的两个方向的变化。
\end{solution}

\begin{exercise}\label{g1:ex:principal-connection-relative-extension}
	设 \(P\to M\) 为光滑主右 \(G\) 丛，\(K\subset M\) 闭，\(U\supset K\) 开，且 \(\pi^{-1}(U)\) 上已给定主联络 \(\omega_U\)。证明存在 \(P\) 上的主联络，在 \(K\) 的某个开邻域上等于 \(\omega_U\)。若两整体主联络均在某个固定开集 \(W\) 上等于指定联络，证明它们可由一条在 \(W\) 上保持该联络不变的光滑联络族连接。
\end{exercise}
\begin{solution}
	由\cref{g1:thm:principal-connection-existence}取任意整体主联络 \(\omega_0\)。由\cref{g1:cor:smooth-cutoff-closed}取 \(\chi:M\to[0,1]\)，使 \(\chi=1\) 于 \(K\) 的某开邻域 \(W_0\subset U\)，且 \(\operatorname{supp}\chi\subset U\)。在 \(\pi^{-1}(U)\) 上，\(\omega_U-\omega_0\) 消去垂直方向并按伴随表示等变。令
	\[
		\omega=\omega_0+(\chi\circ\pi)(\omega_U-\omega_0),
	\]
	第二项在 \(\pi^{-1}(U)\) 外补零。支集包含条件保证补零光滑：在 \(U\) 外每点附近，\(\chi\) 恒为零。权重来自底空间，所以这项仍消去垂直方向并等变；\cref{g1:prop:principal-connection-affine}便保证 \(\omega\) 是主联络。\(\chi=1\) 的地方有 \(\omega=\omega_U\)，满足要求。

	若给定两整体联络 \(\omega_0,\omega_1\)，令 \(\omega_t=(1-t)\omega_0+t\omega_1\)，其中 \(0\leq t\leq1\)。等变条件可直接相加，再现条件给出 \(\omega_t(X^\#)=((1-t)+t)X=X\)，故每个 \(\omega_t\) 都是主联络。系数随 \(t\) 光滑变化；在 \(W\) 上二者原本相同，故整族也保持指定联络不变。
\end{solution}

\begin{exercise}\label{g1:ex:pullback-principal-connection}
	设 \(F:N\to M\) 为无边界光滑流形之间的光滑映射，\(P\to M\) 为主右 \(G\) 丛且带主联络 \(\omega\)。令
	\[
		F^*P=\{(y,p)\in N\times P:F(y)=\pi(p)\},\qquad
		\widehat F(y,p)=p.
	\]
	证明 \(F^*P\to N\) 是主右 \(G\) 丛，\(\widehat F^*\omega\) 是其主联络；不需要 \(F\) 是浸入或淹没。写出拉回联络的局部形式，并说明这一构造与复合光滑映射相容。
\end{exercise}
\begin{solution}
	取 \(P\) 的局部截面 \(s_i:U_i\to P\)。在 \(F^{-1}(U_i)\) 上，
	\[
		(y,g)\longmapsto(y,s_i(F(y))g)
	\]
	给出到 \(F^*P\) 的局部平凡化，逆由 \(P|_{U_i}\) 的平凡化求群坐标得到。其过渡函数为 \(g_{ij}\circ F\)，故光滑；不同底点可由 \(N\) 的开集分离，同一纤维的不同点可在一张乘积图中分离，可数局部图又保证第二可数性。右作用为 \((y,p)g=(y,pg)\)，在这些图中就是群坐标右乘，所以给出主丛。

	记 \(\omega'=\widehat F^*\omega\)。基本场满足 \(\mathrm{d}\widehat F(X^\#_{(y,p)})=X^\#_p\)，因此 \(\omega'(X^\#)=X\)。又因 \(\widehat F\circ R_g=R_g\circ\widehat F\)，有
	\[
		R_g^*\omega'=\widehat F^*(R_g^*\omega)
		=\widehat F^*(\operatorname{Ad}_{g^{-1}}\omega)
		=\operatorname{Ad}_{g^{-1}}\omega'.
	\]
	所以它是主联络。新局部截面为 \(s_i'(y)=(y,s_i(F(y)))\)，故局部形式为
	\[
		(s_i')^*\omega'=F^*(s_i^*\omega)=F^*A_i.
	\]
	以上构造只用局部乘积图与微分的拉回，没有对 \(\mathrm{d}F\) 的秩提出条件。

	若还有 \(H:L\to N\)，则自然映射 \(H^*(F^*P)\to(F\circ H)^*P\) 把 \((z,(H(z),p))\) 送到 \((z,p)\)，在局部群坐标中为恒等，故为主丛同构。两边联络的局部形式分别是 \(H^*(F^*A_i)\) 与 \((F\circ H)^*A_i\)，由链式法则相等，因此自然同构保持联络。
\end{solution}

\begin{exercise}\label{g1:ex:circle-weight-connection-detection}
	设 \(P\to M\) 为主 \(\operatorname{U}(1)\) 丛，\(m\in\mathbb Z\)，取表示 \(\rho_m(z)v=z^m v\) 及关联复线丛 \(L_m=P\times_{\rho_m}\mathbb C\)。
	\begin{enumerate}
		\item 设主联络在局部截面中的形式为 \(A_i\in\Omega^1(U_i;i\mathbb R)\)，求 \(L_m\) 的局部联络形式。
		\item 证明当 \(m\ne0\) 时，两个主联络若在 \(L_m\) 上诱导相同联络，就已经相同；当 \(m=0\) 时，每个主联络都诱导标准平凡线丛上的普通微分。说明这与 \(\rho_m\) 在 \(|m|>1\) 时不忠实并不矛盾。
	\end{enumerate}
\end{exercise}
\begin{solution}
	1. 对 \(X=i a\in i\mathbb R\)，有 \(\rho_m(e^{tX})=e^{tmX}\)，故 \((\rho_m)_*(X)=mX\)。由\eqref{g1:eq:associated-local-derivative}，局部形式为 \(mA_i\)，而对局部分量 \(f_i\)，协变微分为 \(\mathrm{d}f_i+mA_i f_i\)。

	2. 设两主联络的局部形式为 \(A_i,B_i\)。诱导联络相同意味着 \(mA_i=mB_i\)。若 \(m\ne0\)，得到 \(A_i=B_i\)，再由\cref{g1:thm:principal-connection-local-data}的唯一性得到两主联络相同。若 \(m=0\)，表示平凡，\([p,v]\mapsto(\pi(p),v)\) 给出 \(L_0\cong M\times\mathbb C\) 的自然平凡化，局部联络形式处处为零，因此诱导联络就是 \(\mathrm{d}\)。

	当 \(|m|>1\) 时，\(\rho_m\) 的核是 \(|m|\) 个单位根，故它不是忠实群表示；但这个核离散，其 Lie 代数为零，微分 \(X\mapsto mX\) 仍单射。比较同一个主丛上的两主联络，只须判断其 Lie 代数值一形式之差，因此这里仍能全部检出。群表示有非平凡离散核，与微分是否丢失联络方向，是不同的问题。
\end{solution}

\begin{exercise}\label{g1:ex:heisenberg-principal-rectangle}
	令 \(G\) 为实三阶上三角、对角线全为 \(1\) 的矩阵 Lie 群，其 Lie 代数包含 \(X=E_{12}\)、\(Y=E_{23}\)、\(Z=E_{13}\)。在 \(\mathbb R^2\times G\) 上以 \(A=X\,\mathrm{d}x+Y\,\mathrm{d}y\) 定义主联络。计算曲率，并求从 \((0,0,e)\) 出发、沿矩形 \([0,s]\times[0,t]\) 正向边界提升的返回群元素，其中 \(s,t>0\)。说明这里的精确面积公式为何不与一般结构群非交换的情形的限制矛盾。
\end{exercise}
\begin{solution}
	矩阵单位满足 \(XY=Z\)、\(YX=0\)，而 \(Z\) 与 \(X,Y\) 的乘积都为零。因此 \([X,Y]=Z\)，并且 \(Z\) 在这个 Lie 代数中居中。\(A\) 的系数为常值，所以
	\[
		F=\mathrm{d}A+A\wedge A=Z\,\mathrm{d}x\wedge\mathrm{d}y.
	\]
	四条边分别给出左乘 \(e^{-sX}\)、\(e^{-tY}\)、\(e^{sX}\)、\(e^{tY}\)。按实际作用次序，返回元素为
	\[
		h=e^{tY}e^{sX}e^{-tY}e^{-sX}.
	\]
	因为 \(X^2=Y^2=Z^2=0\)，这些指数只有常数项与一次项。直接相乘得到
	\[
		(I+tY)(I+sX)(I-tY)(I-sX)=I-stZ=\exp(-stZ).
	\]
	所以该公式对任意正的 \(s,t\) 精确成立。这里原联络系数不交换，但它们的交换子 \(Z\) 与全部系数交换，更高交换子为零；同时曲率为常值。这些特殊条件消去了更高阶修正，不能由这个例子推出任意结构群非交换的联络都满足同一精确公式。
\end{solution}

\begin{exercise}\label{g1:ex:circle-area-principal-holonomy}
	在 \(\mathbb R^2\times\operatorname{U}(1)\) 上取主联络的局部形式
	\[
		A=\frac{\mathrm i\lambda}{2}(x\,\mathrm{d}y-y\,\mathrm{d}x),\qquad \lambda\in\mathbb R.
	\]
	计算曲率，以及沿半径 \(r>0\) 的正向圆周从 \((r,0,1)\) 出发的返回元素。判断哪些 \(\lambda\) 使联络平坦，并说明返回元素为 \(1\) 的单个圆周能否检测平坦性。
\end{exercise}
\begin{solution}
	\(\operatorname{U}(1)\) 为交换群，故括号项为零，而
	\[
		F=\mathrm{d}A=\mathrm i\lambda\,\mathrm{d}x\wedge\mathrm{d}y.
	\]
	以 \(x=r\cos\theta\)、\(y=r\sin\theta\)、\(0\le\theta\le2\pi\) 参数化，有 \(x\,\mathrm{d}y-y\,\mathrm{d}x=r^2\,\mathrm{d}\theta\)。水平方程为 \(g'=-(\mathrm i\lambda r^2/2)g\)，初值 \(g(0)=1\)，因而
	\[
		g(\theta)=\exp(-\mathrm i\lambda r^2\theta/2),\qquad
		h=\exp(-\mathrm i\lambda\pi r^2).
	\]
	联络平坦当且仅当 \(\lambda=0\)。但只要 \(\lambda r^2\in2\mathbb Z\)，该圆周的返回元素就是 \(1\)；例如 \(\lambda=2/r^2\) 时曲率仍处处非零。单个有限大小的环路只给出一个群元素，指数映射的周期可以使非零面积贡献返回单位元。与此不同，所有充分小圆周的返回都为单位元，会由 \(h=1-\mathrm i\lambda\pi r^2+O(r^4)\) 迫使 \(\lambda=0\)。
\end{solution}

\begin{exercise}\label{g1:ex:associated-curvature-detection}
	设 \(\rho:G\to\operatorname{GL}(V)\) 为有限维实 Lie 群的光滑表示。证明：若 \(\rho_*:\mathfrak g\to\operatorname{End}(V)\) 单射，则关联丛联络平坦当且仅当原主联络平坦。用 \(G=(\mathbb R^2,+)\)、\(\rho(a,b)=e^a\in\operatorname{GL}(1,\mathbb R)\) 构造反例，说明不加这个条件时，关联丛可以看不见非零主曲率。再说明检测曲率并不要求 \(\rho\) 本身单射。
\end{exercise}
\begin{solution}
	在任意局部截面中，关联曲率为 \(\rho_*(F)\)。若 \(F=0\)，显然关联曲率也为零；若 \(\rho_*\) 单射，\(\rho_*(F(v,w))=0\) 又推出每个 \(F(v,w)=0\)，所以反向也成立。各局部截面覆盖底空间，故这个结论是整体的。

	对所给 \(\mathbb R^2\) 表示，有 \(\rho_*(a,b)=a\)。在底空间 \(\mathbb R^2\) 的平凡主丛上取 \(A=(0,x\,\mathrm{d}y)\)，其中群坐标与底空间坐标是不同的两组变量。群为交换群，所以
	\[
		F=(0,\mathrm{d}x\wedge\mathrm{d}y)\ne0,
	\]
	但关联联络形式 \(\rho_*(A)=0\)，故关联丛联络平坦。这是表示在第二个 Lie 代数方向上完全消失所造成的。

	最后，取 \(G=\operatorname{U}(1)\)，以 \(\rho(z)v=z^m v\) 作用于 \(V=\mathbb C\cong\mathbb R^2\)，其中整数 \(m\ge2\)。它的核包含全部 \(m\) 次单位根，所以群表示不单射；但微分把 \(\mathrm i a\) 送到复乘法 \(\mathrm i m a\)，因而单射。离散核不影响 Lie 代数层面的曲率检测，尽管它可能影响整体环路返回的检测。
\end{solution}

\begin{exercise}\label{g1:ex:thin-homotopy-principal-transport}
	设 \(P\to M\) 为带主联络的光滑右主 \(G\) 丛，不假设联络平坦。设 \(\Gamma:[0,1]^2\to M\) 光滑到边界，满足 \(\Gamma(s,0)=x\)、\(\Gamma(s,1)=y\)，并且 \(\mathrm{d}\Gamma\) 在每点的秩至多为 \(1\)。证明所有路径 \(t\mapsto\Gamma(s,t)\) 的主丛运输相同。解释这里为何不需要曲率为零，以及为何这并未证明一般定端点同伦都保持运输。
\end{exercise}
\begin{solution}
	固定 \(p\in P_x\)，令 \(q(s,t)\) 是每条路径从 \(p\) 出发的水平提升。初值曲线常值，所以变分公式给出
	\[
		\omega(q_s(s,1))=\int_0^1\Omega(q_t,q_s)\,\mathrm{d}t.
	\]
	在每个点，\(\Gamma_t,\Gamma_s\) 线性相关。曲率水平，而 \(\mathrm{d}\pi:H\to \mathrm{T}M\) 为同构，故它在 \((q_t,q_s)\) 上的取值只取决于这两个底空间输入的水平提升。两个输入线性相关时，交替二形式取值为零。因此积分右边为零。

	终点 \(\Gamma(s,1)=y\) 固定，故 \(q_s(s,1)\) 是竖直向量；主联络在每个竖直空间上与 \(\mathfrak g\) 同构，所以 \(\omega(q_s)=0\) 迫使 \(q_s(s,1)=0\)。终点曲线因而常值，所有运输把 \(p\) 送到同一点。任意其他初值都唯一写成 \(p g\)，运输的等变性进一步说明全部纤维映射相同。

	这个条件通常称为薄同伦：路径族没有扫出二维切向面积，因而任意曲率二形式在该参数族上的拉回都为零。一般同伦允许微分秩为 \(2\)，这时曲率积分可以非零；本题只消除了当前参数族能够测到的曲率，并没有假设或推出底空间上的曲率消失。
\end{solution}

\begin{exercise}\label{g1:ex:infinitesimal-principal-gauge}
	设 \(P\to M\) 为有限维实 Lie 群 \(G\) 的光滑右主丛，\(\xi\) 为 \(\operatorname{ad}(P)\) 的光滑截面。把它写成等变函数 \(\xi:P\to\mathfrak g\)，令 \(u_t(p)=\exp(t\xi(p))\)。证明 \(u_t\) 给出一族规范变换，并求固定局部截面中联络及曲率的一阶变化。说明这些局部变化为何能拼成整体对象。
\end{exercise}
\begin{solution}
	伴随丛截面的等变条件为 \(\xi(pg)=\operatorname{Ad}_{g^{-1}}\xi(p)\)。指数映射与共轭相容，所以
	\[
		u_t(pg)=\exp(t\operatorname{Ad}_{g^{-1}}\xi(p))
		=g^{-1}\exp(t\xi(p))g.
	\]
	由\cref{g1:prop:gauge-equivariant-functions}，\(\Phi_{u_t}\) 是规范变换。取固定局部截面 \(s_i\)，记 \(\xi_i=\xi\circ s_i\)、\(A_i=s_i^*\omega\)。局部规范公式为
	\[
		A_i(t)=\operatorname{Ad}_{\exp(-t\xi_i)}A_i
		+(\exp(t\xi_i))^*\theta^L.
	\]
	第一项在 \(t=0\) 的导数为 \(-[\xi_i,A_i]=[A_i,\xi_i]\)。第二项的导数为 \(\mathrm{d}\xi_i\)：在群单位元附近取坐标，\(\exp(t\xi_i(x))=e+t\xi_i(x)+O(t^2)\)，而 \(\theta^L_e=\operatorname{id}\)，故对任意底空间切向量 \(v\)，该项为 \(t\,\mathrm{d}\xi_i(v)+O(t^2)\)。这里是坐标中的一阶展开，不要求 \(G\) 为矩阵群。因此
	\[
		\left.\frac{\mathrm{d}}{\mathrm{d}t}\right|_{0}A_i(t)
		=\mathrm{d}\xi_i+[A_i,\xi_i]=\mathrm{D}_{A_i}\xi_i.
	\]
	同理，由 \(F_i(t)=\operatorname{Ad}_{\exp(-t\xi_i)}F_i\)，有
	\[
		\left.\frac{\mathrm{d}}{\mathrm{d}t}\right|_{0}F_i(t)=[F_i,\xi_i].
	\]
	在交叠处，\(\xi_j=\operatorname{Ad}_{g_{ij}^{-1}}\xi_i\)。\cref{g1:cor:adjoint-connection-local}说明 \(\mathrm{D}_A\xi\) 是伴随丛值一形式，局部表达按同一伴随规则变化；曲率与括号的相容性也使 \([F_i,\xi_i]\) 拼成伴随丛值二形式。因而这两项分别是整体联络空间中的切向变化与曲率的变化。特别地，平坦联络沿规范轨道仍平坦。
\end{solution}

\begin{exercise}\label{g1:ex:punctured-plane-flat-principal}
	设 \(p_j=(a_j,b_j)\)、\(1\le j\le k\)，是平面上互异的点，\(k\ge1\)，令 \(M=\mathbb R^2\setminus\{p_1,\ldots,p_k\}\)。定义
	\[
		\alpha_j=\frac{(x-a_j)\,\mathrm{d}y-(y-b_j)\,\mathrm{d}x}
		{(x-a_j)^2+(y-b_j)^2}.
	\]
	在平凡主 \(\operatorname{U}(1)\) 丛上取 \(A=\mathrm{i}\sum_j c_j\alpha_j\)，\(c_j\in\mathbb R\)。求曲率和绕各缺失点正向一周的单值。证明这族联络实现 \(M\) 上全部平坦主 \(\operatorname{U}(1)\) 丛的规范类，并给出参数之间的等价关系。
\end{exercise}
\begin{solution}
	在任何可以选取相对于 \(p_j\) 的角度 \(\theta_j\) 的小邻域内，\(\alpha_j=\mathrm{d}\theta_j\)，所以 \(\mathrm{d}\alpha_j=0\)。\(\operatorname{U}(1)\) 为交换群，故 \(F=\mathrm{d}A=0\)。取从同一基点到 \(p_j\) 附近的小圆周的道路，绕该圆周正向一周后原路返回，记所得环路为 \(\gamma_j\)。来回连接道路的积分相消，小圆周上的积分满足
	\[
		\int_{\gamma_j}\alpha_\ell=2\pi\delta_{j\ell}.
	\]
	其中 \(j=\ell\) 时用圆周角度直接计算；\(j\ne\ell\) 时，\(\alpha_\ell\) 在该小圆周围成的整个圆盘上光滑且闭，由 Stokes 定理积分为零。因此实际返回元素为 \(h_{\gamma_j}=e^{-2\pi\mathrm{i}c_j}\)，本章的左单值为
	\[
		\rho([\gamma_j])=e^{2\pi\mathrm{i}c_j}.
	\]
	\cref{g1:prop:punctured-euclidean}证明 \(\pi_1(M)\) 是以这些环路为生成元的自由群 \(F_k\)。向 \(\operatorname{U}(1)\) 的同态可在自由生成元上任意指定 \(k\) 个相位，而交换目标群中的共轭不改变同态。上述公式实现全部相位，故\cref{g1:thm:flat-principal-classification}说明全部平坦主丛都同构于这族联络中的一员，特别其底层光滑主丛均平凡。规范等价的充要条件是
	\[
		(c_1',\ldots,c_k')-(c_1,\ldots,c_k)\in\mathbb Z^k,
	\]
	所以规范类集合为 \((\mathbb R/\mathbb Z)^k\)。也可以直接实现整数平移：函数
	\[
		u(x,y)=\prod_{j=1}^k
		\left(\frac{x-a_j+\mathrm{i}(y-b_j)}{\sqrt{(x-a_j)^2+(y-b_j)^2}}\right)^{m_j}
	\]
	在 \(M\) 上光滑，且 \(u^{-1}\mathrm{d}u=\mathrm{i}\sum_jm_j\alpha_j\)。这核对了局部规范公式与单值分类中的等价关系。自由群中的非交换信息在这个以交换群为目标的表示下被消去，但缺失点补空间的基本群本身仍是 \(F_k\)，不是预先被改成 \(\mathbb Z^k\)。
\end{solution}

\begin{exercise}\label{g1:ex:so3-gauge-stabilizers}
	设圆周上的平坦主 \(\operatorname{SO}(3)\) 丛的单值生成元是一绕固定轴旋转 \(\theta\) 的矩阵，\(0\le\theta\le\pi\)。求其规范稳定子，分别处理 \(\theta=0\)、\(0<\theta<\pi\)、\(\theta=\pi\)。说明为什么这些稳定子可以发生跳变，而固定纤维基点的规范子群仍自由作用。
\end{exercise}
\begin{solution}
	由\cref{g1:thm:gauge-stabilizer-centralizer}，稳定子同构于单值像的中心化子。实际环路返回是该旋转的逆，但一个元素与其逆具有同一中心化子。选定轴为第三坐标轴，旋转写成 \(R=\operatorname{diag}(R_\theta,1)\)。若 \(\theta=0\)，则 \(R=I\)，其中心化子为整个 \(\operatorname{SO}(3)\)。

	若 \(0<\theta\le\pi\)，\(R\) 的实特征值 \(1\) 的特征空间恰为旋转轴。任何与 \(R\) 交换的正交矩阵都保持这条直线，也保持其正交平面，故写成 \(Q=\operatorname{diag}(B,\varepsilon)\)，其中 \(B\in\operatorname{O}(2)\)、\(\varepsilon=\pm1\)，并有 \(\varepsilon=\det B\)。交换条件归结为 \(BR_\theta=R_\theta B\)。

	当 \(0<\theta<\pi\) 时，写 \(R_\theta=\cos\theta\,I+\sin\theta\,J\)，其中 \(J=\left(\begin{smallmatrix}0&-1\\1&0\end{smallmatrix}\right)\)。因为 \(\sin\theta\ne0\)，条件等价于 \(BJ=JB\)。设 \(B=\left(\begin{smallmatrix}a&b\\c&d\end{smallmatrix}\right)\)，直接比较两边得到 \(d=a\)、\(c=-b\)；再由正交性得 \(a^2+b^2=1\)，所以 \(B\in\operatorname{SO}(2)\)。此时 \(\varepsilon=1\)，中心化子是绕原轴的全部旋转，同构于 \(\operatorname{SO}(2)\)。

	当 \(\theta=\pi\) 时，\(R_\pi=-I\) 在垂直平面上为标量，故任意 \(B\in\operatorname{O}(2)\) 都与它交换。中心化子为
	\[
		\{\operatorname{diag}(B,\det B):B\in\operatorname{O}(2)\}
		\cong\operatorname{O}(2).
	\]
	其中行列式为 \(-1\) 的平面变换还反转旋转轴，这部分在一般角度时不再交换。因此稳定子不仅在恒等单值处增大，在半周旋转处也多出一个连通分支。固定 \(p_0\) 的规范子群则要求稳定子对应元素 \(u(p_0)=e\)，由同一定理只能是恒等规范变换，所以其作用一直自由；这不与全规范群的稳定子跳变矛盾。
\end{solution}

\begin{exercise}\label{g1:ex:hopf-powers-flat-obstruction}
	在 Hopf 丛的两张射影图上，令 \(g_{01}(w)=|w|/w\)，并以 \(g_{01}^m\) 为过渡函数构造主 \(\operatorname{U}(1)\) 丛 \(P_m\to\mathbb{CP}^1\)，其中 \(m\in\mathbb Z\)。证明原 Hopf 联络的局部形式 \(A_0,A_1\) 乘以 \(m\) 后给出 \(P_m\) 上的联络，求其曲率积分，并判断 \(P_m\) 何时平凡、何时容许平坦联络。
\end{exercise}
\begin{solution}
	连同 \(g_{10}^m=(g_{01}^m)^{-1}\)、\(g_{ii}^m=1\)，这组函数满足余循环条件，故\cref{g1:thm:principal-cocycle-gluing}给出 \(P_m\)。原来的局部联络满足 \(A_1=A_0+g_{01}^{-1}\mathrm{d}g_{01}\)。对整数 \(m\)，包括负整数，由乘积求导及交换性可得
	\[
		(g_{01}^m)^{-1}\mathrm{d}(g_{01}^m)
		=m\,g_{01}^{-1}\mathrm{d}g_{01}.
	\]
	于是 \(mA_1=mA_0+(g_{01}^m)^{-1}\mathrm{d}(g_{01}^m)\)，满足主联络的粘合公式。其曲率是 \(F_m=mF\)，按正文的复定向有
	\[
		\frac1{2\pi\mathrm{i}}\int_{\mathbb{CP}^1}F_m=m.
	\]
	若 \(m\ne0\)，全局截面会使 \(F_m\) 成为整体恰当二形式，Stokes 定理迫使积分为零，矛盾，故 \(P_m\) 不平凡。对任意另一联络，\(\operatorname{U}(1)\) 的伴随表示平凡，联络之差下降为底空间上的整体 \(\mathrm{i}\mathbb R\) 值一形式 \(\beta\)，曲率之差为 \(\mathrm{d}\beta\)。积分因而仍是 \(2\pi\mathrm{i}m\)，所以没有平坦联络。

	若 \(m=0\)，过渡函数全部为 \(1\)，主丛为乘积，局部形式 \(mA_i=0\) 粘成乘积平坦联络。因此平凡性和容许平坦联络在这族中都恰好等价于 \(m=0\)。这里断言的是存在平坦联络，不是说平凡丛上的每个联络都平坦；后一说法在二维底空间上并不成立。
\end{solution}
